% !TEX root = 抽象代数前三章全部习题详解.tex


\clearpage\section{第3章：环}


\noindent（下列各题中出现的环，除相反的声明外，都假定是有单位元素的环。）

原题逐页核对聂灵沼、丁石孙《代数学引论》第三版，书页131--136（本地扫描PDF第141--146页）。各题保留原来的题意、子问顺序、定义和提示，只转换为TeX排版。原题的约定问题与可核实的错误单列在解答中，题面不作替换。下文证明为本册解答，使用原笔记的记号与方法；教材仅给题目时，不声称书中给出了这些证明。


\subsection{习题1：由\texorpdfstring{$1-ab$}{1-ab}可逆推出\texorpdfstring{$1-ba$}{1-ba}可逆}

\label{ex:3:1}

\paragraph{题目}

证明在环$R$内，若$1-ab$有逆，则$1-ba$也有逆。

\paragraph{考点与思路}

在非交换环中，不能将 $ab$ 与 $ba$ 直接互换。先记 $c=(1-ab)^{-1}$，再尝试把原来的逆元嵌入 $bca$，构造 $1-ba$ 的逆元。必须分别验证左右两个乘积。

\paragraph{所用结论与原文依据}

\noteref{ringlaws}；\noteref{unit}。原题：《代数学引论》第三版，书页131，扫描PDF第141页。

% 原题逐页视觉核对；扫描PDF第141页；本题证明为本册解答。

\begin{proof}

由 $c=(1-ab)^{-1}$ 的定义，
\[
(1-ab)c=1,\qquad c(1-ab)=1.
\]
把两式分别展开，得到
\[
c-abc=1,\qquad c-cab=1.
\]
令 $d=1+bca$。先计算 $(1-ba)d$，使用结合律把含 $b$ 和 $a$ 的项提取出来：
\begin{align*}
(1-ba)(1+bca)
  &=1+bca-ba-babca\\
  &=1+b(c-1-abc)a\\
  &=1.
\end{align*}
最后一步用了 $c-abc=1$。

另一方向需要使用另一条逆元等式：
\begin{align*}
(1+bca)(1-ba)
  &=1-ba+bca-bcaba\\
  &=1+b(c-1-cab)a\\
  &=1,
\end{align*}
其中最后一步用了 $c-cab=1$。因此 $d$ 是 $1-ba$ 的双侧逆元，且
\[
\boxed{(1-ba)^{-1}=1+b(1-ab)^{-1}a.}
\]

\end{proof}

\paragraph{复习提示}

记住逆元的构造次序是 $1+bca$。证明中没有交换任意两个因子，因此适用于非交换幺环。将 $a,b$ 对换也得到逆命题，故 $1-ab$ 可逆当且仅当 $1-ba$ 可逆。



\subsection{习题2：右逆不唯一、左零因子与非单位的等价性}

\label{ex:3:2}

\paragraph{题目}

设在环$R$中元素$u$有右逆，证明下列三条等价：
\begin{enumerate}[label=(\roman*)]
\item $u$有多于一个的右逆；
\item $u$是一个左零因子；
\item $u$不是单位。
\end{enumerate}

\paragraph{考点与思路}

先固定一个右逆 $v$，即 $uv=1$。两个右逆相减后落入方程 $ut=0$ 的解集中；反过来，给 $v$ 加上这样的 $t$ 就能产生新的右逆。条件 $uv=1$ 还排除了非零元素在 $u$ 左侧乘出零的可能。

\paragraph{所用结论与原文依据}

\noteref{ringlaws}；\noteref{unit}。原题：《代数学引论》第三版，书页131，扫描PDF第141页。

% 原题逐页视觉核对；扫描PDF第141页；本题证明为本册解答。

\begin{proof}

固定 $v\in R$ 满足 $uv=1$。若 $R$ 是只有一个元素的零环，则 $u$ 是单位，只有一个右逆且不是零因子，三个条件均不成立，等价性仍成立。以下设 $1\ne0$，于是 $u\ne0$。

\begin{enumerate}\item 证明条件(1)与条件(2)等价。
若另有 $w\ne v$ 满足 $uw=1$，则 $w-v\ne0$，且
\[
u(w-v)=uw-uv=1-1=0.
\]
因此 $u$ 是左零因子，特别是零因子。

反过来，先注意到：对于任意 $x\in R$，若 $xu=0$，则
\[
x=x\cdot1=x(uv)=(xu)v=0.
\]
所以不存在非零 $x$ 使 $xu=0$。若 $u$ 是零因子，必有非零 $t$ 满足 $ut=0$。于是
\[
u(v+t)=uv+ut=1,
\]
而 $v+t\ne v$，故 $u$ 至少有两个右逆。

\item 证明条件(1)与条件(3)等价。
若 $u$ 是单位，任取右逆 $w$，由 $uw=1$ 得
\[
w=u^{-1}(uw)=u^{-1}.
\]
因而右逆唯一。所以右逆不唯一必推出 $u$ 不是单位。

另一方面，设 $u$ 的右逆唯一。注意
\[
u(vu-1)=(uv)u-u=u-u=0.
\]
因此 $v+(vu-1)$ 也是右逆，因为
\[
u\bigl(v+vu-1\bigr)=uv+u(vu-1)=1.
\]
由右逆的唯一性，$v+vu-1=v$，在加法群中消去 $v$，得到 $vu=1$。结合 $uv=1$，可知 $u$ 是单位。

取这条结论的逆否命题，$u$ 不是单位便推出其右逆不唯一。至此三个条件两两等价。
\end{enumerate}

\end{proof}

\paragraph{复习提示}

这里的前提“已经有一个右逆”不可遗漏。一般幺环中“不是单位”未必是零因子，例如整数环中的 $2$。本题还证明：有右逆的非单位只能以 $ut=0$ 的形式成为零因子，不能有非零 $x$ 满足 $xu=0$。



\subsection{习题3：右逆不唯一便有无穷多个右逆}

\label{ex:3:3}

\paragraph{题目}

在环$R$内，若元素$u$有多于一个的右逆，证明它有无穷多个右逆。

\paragraph{考点与思路}

由 $uv=1$ 但 $vu\ne1$ 构造非零元素 $e=1-vu$。它满足 $ue=0$、$ev=0$，适合用来给已有的右逆 $v$ 添加不同的修正项。不能仅写 $v+ne$，因为正特征下这些元素可能重复。

\paragraph{所用结论与原文依据}

\noteref{ringlaws}；\noteref{unit}。原题：《代数学引论》第三版，书页131，扫描PDF第141页。

% 原题逐页视觉核对；扫描PDF第141页；本题证明为本册解答。

\begin{proof}

取一个右逆 $v$，使 $uv=1$。由\exref{3}{2}，$u$ 不是单位。若 $vu=1$，则 $v$ 同时是左逆和右逆，$u$ 便是单位，矛盾。因此
\[
e=1-vu\ne0.
\]
由 $uv=1$，有
\begin{align*}
ue&=u(1-vu)=u-(uv)u=0,\\
ev&=(1-vu)v=v-v(uv)=0.
\end{align*}

对每个整数 $k\ge0$，令
\[
x_k=v+eu^k,
\]
这里 $u^0=1$。因为 $ue=0$，
\[
ux_k=uv+(ue)u^k=1,
\]
所以每个 $x_k$ 都是 $u$ 的右逆。

下面证明它们两两不同。先由 $uv=1$ 归纳得到，对每个整数 $m\ge0$，
\[
u^mv^m=1.
\]
归纳步骤为
\[
u^{m+1}v^{m+1}=u^m(uv)v^m=u^mv^m=1.
\]
于是当 $0\le i<j$ 时，
\[
u^iv^j=(u^iv^i)v^{j-i}=v^{j-i}.
\]

假设 $x_i=x_j$，其中 $i<j$。消去 $v$ 后，$eu^i=eu^j$。两边右乘 $v^j$，得到
\[
ev^{j-i}=e.
\]
但 $j-i\ge1$ 且 $ev=0$，所以
\[
ev^{j-i}=(ev)v^{j-i-1}=0,
\]
从而 $e=0$，与构造矛盾。因此 $x_0,x_1,x_2,\ldots$ 两两不同，$u$ 有无穷多个右逆。

\end{proof}

\paragraph{复习提示}

证明的核心是“构造右逆”和“证明构造出的元素互异”两个步骤。等式 $u^mv^m=1$ 来自逐对消去中间的 $uv$，绝不意味着 $u,v$ 可以交换。本题也说明有限幺环中，有右逆的元素必为单位。



\subsection{习题4：幂零元全体是否为理想}

\label{ex:3:4}

\paragraph{题目}

已知一个交换环$R$的全部零幂元构成$R$的一个理想，在非交换环内，上述事实是否成立？

\paragraph{考点与思路}

交换情形展开二项式，并分别检查减法和吸收性；非交换情形保留两个矩阵单位的反例。

\paragraph{所用结论与原文依据}

\noteref{ringlaws}；\noteref{ideal}；\noteref{c3nil}。原题：《代数学引论》第三版，书页131，扫描PDF第141页。

% 原题逐页视觉核对；扫描PDF第141页；本题证明为本册解答。

\begin{proof}

令$N=\{a\in R:a^m=0\text{ 对某个 }m\ge1\}$。因为$0\in N$，它非空。
任取$a,b\in N$，选$m,n\ge1$使$a^m=b^n=0$。因$R$交换，二项式展开的每一项形如$a^i(-b)^j$，其中$i+j=m+n-1$。若$i<m$且$j<n$，则$i+j\le m+n-2$，矛盾。因此每一项均含$a^m$或$b^n$，从而
\[
(a-b)^{m+n-1}=0.
\]
这说明$a-b\in N$。

任取$r\in R$，交换性给出$(ra)^m=r^ma^m=0$，故$ra=ar\in N$。所以$N$是理想。

在非交换环$M_2(\Z)$中，取$A=E_{12}$、$B=E_{21}$。有
\[
A^2=B^2=0,\qquad (A+B)^2=E.
\]
于是$(A+B)^{2k}=E\ne0$，$(A+B)^{2k+1}=A+B\ne0$，所以$A+B$不幂零。幂零元全体甚至不对加法封闭，因而一般不是理想。

\end{proof}

\paragraph{复习提示}

二项式展开中将同类乘积合并使用了交换性；不能把该证明直接移用于非交换环。



\subsection{习题5：素特征的幂公式与幺环的特征}

\label{ex:3:5}

\paragraph{题目}

设$F$为一个特征$p$的域，$p$为素数，证明
\[
(a+b)^p=a^p+b^p,\qquad\text{对所有的 }a,b\in F.
\]
域的特征概念可以推广到幺环上去。如果幺环$R$的单位元素$e$生成的加法群$G=\{ne\mid n\in\Z\}$是一个无限群，则$R$叫做特征$0$的环；若$G$是一个有限群，令$k=|G|$，则$k$叫做环$R$的特征。证明：若$R$为整环，则$R$的特征为$0$或为一个素数，而且若$R$的特征为素数$p$，则上面等式对$R$也成立。

\paragraph{考点与思路}

先证明中间二项式系数被素数$p$整除，再用单位元素的加法阶讨论整环的特征。

\paragraph{所用结论与原文依据}

\noteref{ringlaws}；\noteref{unit}；\noteref{zerodiv}；\noteref{field}。原题：《代数学引论》第三版，书页131--132，扫描PDF第141页起。

% 原题逐页视觉核对；扫描PDF第141页；本题证明为本册解答。

\begin{proof}

\begin{enumerate}
\item 对$1\le k\le p-1$，整数等式
\[
k\binom pk=p\binom{p-1}{k-1}
\]
和$p\nmid k$说明$p\mid\binom pk$。因此这些二项式系数在特征$p$的域$F$中均为零。由$F$交换，
\[
(a+b)^p=a^p+\sum_{k=1}^{p-1}\binom pk a^kb^{p-k}+b^p=a^p+b^p.
\]
\item 现设$R$为整环，单位元记为$e$。若$e$在加法下的阶无限，按题中定义，特征为$0$。若其阶为正整数$k$，因为整环$e\ne0$，有$k>1$。

假设$k$不是素数，写$k=mn$，其中$1<m,n<k$。$k$的最小性给出$me\ne0$、$ne\ne0$，但
\[
(me)(ne)=(mn)e=ke=0,
\]
违背整环无零因子。因此$k$为素数。
\item 若此整环$R$的特征为$p$，则$pr=(pe)r=0$对任意$r\in R$成立，因而凡被$p$整除的整数系数都作用为零。整环交换，第一项的二项式证明逐项适用，故$(a+b)^p=a^p+b^p$对$R$中全部$a,b$也成立。
\end{enumerate}

\end{proof}

\paragraph{复习提示}

最后一句仍在“$R$为整环”的条件下讨论；不能把交换性省去。



\subsection{习题6：矩阵环的全部双边理想}

\label{ex:3:6}

\paragraph{题目}

设$M_n(R)$为环$R$上$n\times n$全矩阵环，又设$I$为$R$的一个理想，求证$I$上$n\times n$全矩阵环$M_n(I)$是$M_n(R)$的理想。进一步证明，$M_n(R)$除这一类理想外再无其它理想。

\paragraph{考点与思路}

用矩阵单位取出一个分量并把它搬到任意位置，再重组整个矩阵。

\paragraph{所用结论与原文依据}

\noteref{ideal}；\noteref{unit}。原题：《代数学引论》第三版，书页132，扫描PDF第142页。

% 原题逐页视觉核对；扫描PDF第142页；本题证明为本册解答。

\begin{proof}

这里$R$为幺环，$n\ge1$，理想均为双边理想。若$I\lhd R$，则$M_n(I)$对减法封闭；对$A\in M_n(R)$、$B\in M_n(I)$，乘积$AB,BA$的每个分量都是含有$I$中一个因子的有限和，仍属于$I$。故$M_n(I)\lhd M_n(R)$。

反过来，设$J\lhd M_n(R)$。令
\[
I=\{a\in R:aE_{11}\in J\}.
\]
零矩阵属于$J$，故$0\in I$。若$a,b\in I$、$r\in R$，则
\begin{align*}
(a-b)E_{11}&=aE_{11}-bE_{11}\in J,\\
(ra)E_{11}&=(rE_{11})(aE_{11})\in J,\\
(ar)E_{11}&=(aE_{11})(rE_{11})\in J.
\end{align*}
因此$I\lhd R$。

对$A=(a_{ij})\in J$，矩阵单位给出
\[
E_{1i}AE_{j1}=a_{ij}E_{11}\in J,
\]
所以每个$a_{ij}\in I$，即$J\subseteq M_n(I)$。反之，对$a\in I$及任意$i,j$，
\[
E_{i1}(aE_{11})E_{1j}=aE_{ij}\in J.
\]
将这些矩阵相加，$M_n(I)$中的每个矩阵都属于$J$，故$J=M_n(I)$。由$a\in I\Longleftrightarrow aE_{11}\in J$，$I$也唯一。

\end{proof}

\paragraph{复习提示}

所有分量都能通过矩阵单位单独取出，并能搬到任意位置，所以能恢复且唯一确定系数理想。



\subsection{习题7：理想积的分配律}

\label{ex:3:7}

\paragraph{题目}

设$I,H,N$为环$R$的理想，证明
\[
(H+N)I=HI+NI,\qquad I(H+N)=IH+IN.
\]

\paragraph{考点与思路}

按理想积的有限和定义，逐项展开，再分别证明反向包含。

\paragraph{所用结论与原文依据}

\noteref{c3product}；\noteref{ringlaws}。原题：《代数学引论》第三版，书页132，扫描PDF第142页。

% 原题逐页视觉核对；扫描PDF第142页；本题证明为本册解答。

\begin{proof}

按定义，$HI$是全部有限和$\sum h_ji_j$组成的集合，其中$h_j\in H,i_j\in I$；其余理想积也按同样方式理解。

任取$x\in(H+N)I$，写作$x=\sum_j(h_j+n_j)i_j$。分配律给出
\[
x=\sum_jh_ji_j+\sum_jn_ji_j\in HI+NI.
\]
反过来，因$H,N\subseteq H+N$，每个$h_ji_j$和$n_ki_k$都是$(H+N)I$的生成乘积；该集合对加法封闭，故$HI+NI\subseteq(H+N)I$。于是$(H+N)I=HI+NI$。

对另一等式，任取$y=\sum_ji_j(h_j+n_j)\in I(H+N)$，则
\[
y=\sum_ji_jh_j+\sum_ji_jn_j\in IH+IN.
\]
而$H,N\subseteq H+N$分别给出$IH,IN\subseteq I(H+N)$，故$IH+IN\subseteq I(H+N)$。得到$I(H+N)=IH+IN$。两式均不要求乘法交换。

\end{proof}

\paragraph{复习提示}

理想积的元素是有限个乘积之和，不一定是单个乘积。



\subsection{习题8：满同态下的元素性质与整环性质}

\label{ex:3:8}

\paragraph{题目}

设$\eta:R\to R'$是一个满的环同态而且将$R$的单位元素$1$映到单位元素$1'$。指出下列命题的正确和错误。正确的给以证明，错误的请举反例：
\begin{enumerate}[label=(\roman*)]
\item 若$a\in R$是幂零（幂等）元，则$\eta(a)$也是$R'$的幂零（幂等）元（如果环的元素$e$适合$e^2=e$，则$e$叫做幂等元）；
\item 若$a\in R$是零因子，则$\eta(a)$也是$R'$的零因子；
\item 若$R$为整环，则$\eta(R)=R'$也为整环；
\item 若$\eta(R)=R'$为整环，则$R$也是整环；
\item 若$u\in R$为单位，则$\eta(u)$也是$R'$的单位；
\item 若$\eta(u)$是$R'$的单位，则$u$也是$R$的单位。
\end{enumerate}

\paragraph{考点与思路}

同态保留由加法、乘法和 $0,1$ 写出的等式，却未必保留“非零”“不存在非零零乘积”等条件。零因子的定义涉及非零元素，所以不能只把 $ab=0$ 两边取像就认定像仍为零因子。

\paragraph{所用结论与原文依据}

\noteref{ringiso}；\noteref{ideal}。原题：《代数学引论》第三版，书页132，扫描PDF第142页。

% 原题逐页视觉核对；扫描PDF第142页；本题证明为本册解答。

\begin{proof}

\begin{enumerate}[label=(\roman*)]\item 成立。
若 $a$ 是幂零元，取正整数 $m$ 使 $a^m=0$。由同态保乘法，
\[
\eta(a)^m=\eta(a^m)=\eta(0)=0'.
\]
故 $\eta(a)$ 仍为幂零元。若 $e^2=e$，则
\[
\eta(e)^2=\eta(e^2)=\eta(e),
\]
故 $\eta(e)$ 仍为幂等元。两项结论都不要求满射或保幺。

\item 不成立。
考虑映射
\[
\eta:\Z/6\Z\longrightarrow\Z/3\Z,\qquad
\overline{k}_{6}\longmapsto\overline{k}_{3}.
\]
若 $k\equiv\ell\pmod6$，则 $3\mid k-\ell$，所以像不依赖代表元，映射良定义。模运算保留加法和乘法，且 $\eta(\overline1_6)=\overline1_3$；每个 $\overline{k}_3$ 都是 $\overline{k}_6$ 的像，故这是保幺满同态。

在源环中，$\overline4_6\ne0$、$\overline3_6\ne0$，且
\[
\overline4_6\,\overline3_6=\overline{12}_6=\overline0_6,
\]
所以 $\overline4_6$ 是零因子。它的像却是 $\overline1_3$。若 $\overline1_3 y=0$，则 $y=0$，故该像不是零因子。这一反例甚至将零因子映成了非零单位。

\item 不成立。
取自然映射
\[
\pi:\Z\twoheadrightarrow\Z/6\Z,\qquad k\longmapsto\overline{k}_6.
\]
该映射保加法、乘法和单位元，并且满射。源环 $\Z$ 是整环，但目标中
\[
\overline2_6\ne0,\qquad \overline3_6\ne0,\qquad
\overline2_6\,\overline3_6=0,
\]
所以 $\Z/6\Z$ 不是整环。

\item 不成立。
取映射
\[
\rho:\Z/6\Z\twoheadrightarrow\Z/2\Z,\qquad
\overline{k}_{6}\longmapsto\overline{k}_{2}.
\]
因为 $2\mid6$，同余模 $6$ 蕴含同余模 $2$，映射良定义；模运算说明它保加法、乘法及单位元，每个模 $2$ 的剩余类都有原像，所以满射。

目标 $\Z/2\Z$ 只有两个元素，唯一非零元素是 $\overline1_2$，且 $\overline1_2\,\overline1_2=\overline1_2\ne0$，因此它是整环。源环却有上一项给出的非零零乘积，故不是整环。

\item 成立。
若 $u$ 是单位，令 $v=u^{-1}$，则 $uv=vu=1$。取像得
\begin{align*}
\eta(u)\eta(v)&=\eta(uv)=\eta(1)=1',\\
\eta(v)\eta(u)&=\eta(vu)=\eta(1)=1'.
\end{align*}
因此 $\eta(u)$ 是单位，且
\[
\eta(u)^{-1}=\eta(u^{-1}).
\]
这一结论只需要保幺同态，不需要满射。

\item 不成立。
取自然保幺满同态
\[
\sigma:\Z\twoheadrightarrow\Z/5\Z,\qquad
k\longmapsto\overline{k}_5.
\]
整数 $2$ 不是 $\Z$ 中的单位：若 $2m=1$，则偶数等于奇数，故没有整数解。它的像 $\overline2_5$ 却满足
\[
\overline2_5\,\overline3_5=\overline6_5=\overline1_5,
\]
因而是单位。这说明像可逆不保证原元素可逆。
\end{enumerate}

\end{proof}

\paragraph{复习提示}

六项判断依次为：\textbf{对、错、错、错、对、错}。构造剩余类环之间的映射时，先检验模数之间的整除关系，确保“按原整数取剩余类”的规则良定义。证明零因子的反例时，须同时写明两个因子非零；证明单位的反例时，须给出像的实际逆元。



\subsection{习题9：有理系数四元数构成子体}

\label{ex:3:9}

\paragraph{题目}

在四元数体$H$内，令
\[
H_0=\{a_0+a_1I+a_2J+a_3K\mid a_i\in\Q\}.
\]
证明$H_0$是$H$的子体。

\paragraph{考点与思路}

沿用四元数乘法表检查封闭性，再用共轭除以平方和求逆。

\paragraph{所用结论与原文依据}

\noteref{quaternion}；\noteref{c3conjugate}；\noteref{field}。原题：《代数学引论》第三版，书页132，扫描PDF第142页。

% 原题逐页视觉核对；扫描PDF第142页；本题证明为本册解答。

\begin{proof}

$H_0$含$0,1$，且对减法封闭。将两个有理系数四元数的乘积按$1,I,J,K$展开，由乘法表可知每个新系数都是原系数的整数系数和积，仍为有理数。因此$H_0$是$H$的含幺子环。

任取$0\ne q=a_0+a_1I+a_2J+a_3K\in H_0$，定义$\bar q=a_0-a_1I-a_2J-a_3K$。不同虚基元的交叉项成对抵消，得到
\[
q\bar q=\bar q q=a_0^2+a_1^2+a_2^2+a_3^2=:s.
\]
至少一个$a_i$非零，且这些系数属于$\Q\subset\R$，所以$s>0$。特别地，$s$是非零有理数，故
\[
q^{-1}=s^{-1}\bar q\in H_0.
\]
左右乘$q$均得到$1$，所以每个非零元素都有逆，$H_0$是$H$的子体。这里“体”指除环；$IJ=K$、$JI=-K$说明它不是交换域。

\end{proof}

\paragraph{复习提示}

平方和非零的理由来自系数位于实数中；在任意域上不能直接使用正性。



\subsection{习题10：实四元数体的中心}

\label{ex:3:10}

\paragraph{题目}

决定四元数体$H$的中心。

\paragraph{考点与思路}

先令一般四元数分别与$I,J$交换，以系数唯一性逐次排除虚部。

\paragraph{所用结论与原文依据}

\noteref{quaternion}；\noteref{c3conjugate}。原题：《代数学引论》第三版，书页132，扫描PDF第142页。

% 原题逐页视觉核对；扫描PDF第142页；本题证明为本册解答。

\begin{proof}

环的中心定义为$Z(H)=\{q\in H:qr=rq\text{ 对每个 }r\in H\}$。任取$q=a+bI+cJ+dK\in Z(H)$。由乘法表，
\begin{align*}
qI&=-b+aI+dJ-cK,\\
Iq&=-b+aI-dJ+cK.
\end{align*}
由$qI=Iq$及$1,I,J,K$的实线性无关性，有$2d=2c=0$，故$c=d=0$。

此时$q=a+bI$，又
\[
qJ=aJ+bK,\qquad Jq=aJ-bK.
\]
因$qJ=Jq$，得$2b=0$，故$b=0$。所以$Z(H)\subseteq\R1$。
反过来，任一实标量$a1$与四元数的各基元交换，对分配律展开后的任意四元数也交换。因此
\[
Z(H)=\R1\cong\R.
\]

\end{proof}

\paragraph{复习提示}

必要性只需与两个特定元素交换；充分性仍须说明实标量与所有四元数交换。



\subsection{习题11：与四元数基元乘法交换的元素}

\label{ex:3:11}

\paragraph{题目}

决定四元数体$H$中与$I$乘法交换的元素全体。

\paragraph{考点与思路}

写出一般四元数与$I$相乘的两个结果，以系数相等求出全部交换元。

\paragraph{所用结论与原文依据}

\noteref{quaternion}。原题：《代数学引论》第三版，书页132，扫描PDF第142页。

% 原题逐页视觉核对；扫描PDF第142页；本题证明为本册解答。

\begin{proof}

任取$q=a+bI+cJ+dK\in H$，其中$a,b,c,d\in\R$。由四元数乘法表，
\begin{align*}
Iq&=-b+aI-dJ+cK,\\
qI&=-b+aI+dJ-cK.
\end{align*}
所以
\[
Iq-qI=2cK-2dJ.
\]
$J,K$在实数域上线性无关，故$I q=q I$当且仅当$c=d=0$。因此所求元素全体为
\[
\boxed{\{a+bI:a,b\in\R\}.}
\]
反过来，对任意$a+bI$，两边乘$I$都得到$-b+aI$，所以这些元素全部与$I$交换，分类没有遗漏。

\end{proof}

\paragraph{复习提示}

原题求的是交换元全体，不是它的补集。



\subsection{习题12：共轭不变的四元数子体}

\label{ex:3:12}

\paragraph{题目}

设$S$是四元数体$H$的一个子体，若对所有非零元素$d\in H$都有$dSd^{-1}\subset S$，证明$S=H$或者$S$被包含在$H$的中心内。

\paragraph{考点与思路}

取$x$与$x+1$的两个共轭，比较后迫使它们相等，再以$x+y$处理子体内部的元素。

\paragraph{所用结论与原文依据}

\noteref{field}；\noteref{c3conjugate}；\exref{3}{10}。原题：《代数学引论》第三版，书页132，扫描PDF第142页。

% 原题逐页视觉核对；扫描PDF第142页；本题证明为本册解答。

\begin{proof}

若$S=H$则结论成立。以下设$S\ne H$，固定$x\in H\setminus S$。子体$S$含$0,1$，所以$x\ne0$且$x+1\ne0$。任取$s\in S$，题设保证
\[
a=xsx^{-1}\in S,\qquad b=(x+1)s(x+1)^{-1}\in S.
\]
右乘$x$及$x+1$分别得$xs=ax$及$xs+s=bx+b$。相减得
\[
(a-b)x=b-s.
\]
若$a\ne b$，则$a-b$在子体$S$中可逆，从左侧消去得$x=(a-b)^{-1}(b-s)\in S$，矛盾。故$a=b$，上述等式进而给出$b=s$，于是$xs=sx$。

由于$x$在$H\setminus S$中可任取，子体外的每个元素都与$s$交换。对任意$y\in S$，固定一个$x\notin S$；若$x+y\in S$则$x=(x+y)-y\in S$，矛盾。因此$x+y\notin S$，它也与$s$交换。于是
\[
ys=(x+y)s-xs=s(x+y)-sx=sy.
\]
所以$s$与$S$内外的所有元素均交换，即$s\in Z(H)$。$s$任意，故$S\subseteq Z(H)=\R1$。

\end{proof}

\paragraph{复习提示}

最后以$x+y$相减，才能从子体外的交换性推出子体内部的交换性。



\subsection{习题13：互素理想与理想积的包含}

\label{ex:3:13}

\paragraph{题目}

设$I,H,N$为环$R$的理想。若$HN\subset I$而且$I$与$H$互素，证明$N\subset I$。

\paragraph{考点与思路}

由$I+H=R$把$1$分成两项，再乘以$N$中的任意元素。

\paragraph{所用结论与原文依据}

\noteref{ideal}；\noteref{c3product}。原题：《代数学引论》第三版，书页132，扫描PDF第142页。

% 原题逐页视觉核对；扫描PDF第142页；本题证明为本册解答。

\begin{proof}

这里$R$为幺环。由$I+H=R$，可取$i\in I,h\in H$使$1=i+h$。任取$n\in N$，有
\[
n=1n=(i+h)n=in+hn.
\]
因$I$为理想，$in\in I$；又$h\in H,n\in N$，所以$hn\in HN\subseteq I$。故$n\in I$。$n$任意，得到$N\subseteq I$。

\end{proof}

\paragraph{复习提示}

利用原书共同约定的幺元，先将$1$分解成$I$与$H$中的两项。



\subsection{习题14：互素理想之积包含其交中的理想}

\label{ex:3:14}

\paragraph{题目}

设$I,H,N$是交换环$R$的理想。若$H\supset I,N\supset I$而且$H,N$互素，证明$HN\supset I$。

\paragraph{考点与思路}

由$1=h+n$把$I$中的元素拆成两个属于$HN$的项。

\paragraph{所用结论与原文依据}

\noteref{c3comaximal}；\noteref{c3product}。原题：《代数学引论》第三版，书页132，扫描PDF第142页。

% 原题逐页视觉核对；扫描PDF第142页；本题证明为本册解答。

\begin{proof}

这里$R$为交换幺环。由$H+N=R$，取$h\in H,n\in N$使$h+n=1$。任取$i\in I$。因$I\subseteq H\cap N$，有$i\in H$且$i\in N$，所以
\[
i=i(h+n)=ih+in=hi+in\in HN.
\]
其中$hi\in HN$使用$h\in H,i\in N$，$in\in HN$使用$i\in H,n\in N$。故$I\subseteq HN$。

\end{proof}

\paragraph{复习提示}

对$I$中的元素乘以$1=h+n$，逐项检查属于$HN$的理由。



\subsection{习题15：只有恒等反自同构的矩阵子环}

\label{ex:3:15}

\paragraph{题目}

在域$F$上$n\times n$（$n>1$）全矩阵环$M_n(F)$内寻求一个子环$R$，使得$R$除恒等同构外没有其它反自同构。

\paragraph{考点与思路}

取素子域的标量矩阵环；先证明自同构必须固定幺元，再固定全部素域元素。

\paragraph{所用结论与原文依据}

\noteref{c3anti}；\noteref{unit}；\noteref{field}。原题：《代数学引论》第三版，书页132，扫描PDF第142页。

% 原题逐页视觉核对；扫描PDF第142页；本题证明为本册解答。

\begin{proof}

令$F_0$为$F$的素子域，即$\operatorname{char}F=0$时的$\Q$，或特征为$p$时的$\F_p$，并取
\[
R=\{cE_n:c\in F_0\}\subseteq M_n(F).
\]
标量矩阵的加法、乘法对应于$F_0$的运算，故$R$是交换子环，单位元为$E_n$。

设$\sigma:R\to R$为反自同构。因$R$交换，$\sigma(AB)=\sigma(B)\sigma(A)=\sigma(A)\sigma(B)$，故它也是自同构。即使不预设保幺，满射性也迫使$\sigma(E_n)=E_n$：对任意$B=\sigma(A)$，都有$\sigma(E_n)B=B\sigma(E_n)=B$，再利用单位元唯一性。

加法性说明$\sigma$固定每个整数倍$mE_n$。正特征时这些就是全部$F_0$标量。零特征时，对非零整数$m$，由$\sigma$保乘法，$\sigma((mE_n)^{-1})=(mE_n)^{-1}$，继而固定每个有理数标量。因此两种情形都只有$\sigma=\mathrm{id}_R$。

\end{proof}

\paragraph{复习提示}

题目只要求一个子环，素子域标量环已经满足要求；不必让子环等于整个矩阵环。



\subsection{习题16：整数下三角矩阵环的全部理想}

\label{ex:3:16}

\paragraph{题目}

设$R=\left\{\begin{pmatrix}a&0\\b&c\end{pmatrix}\middle|a,b,c\in\Z\right\}$，试决定$R$的全部理想。

\paragraph{考点与思路}

用对角矩阵单位分离三个位置，再用$E_{21}$比较对应的三个整数理想。

\paragraph{所用结论与原文依据}

\noteref{ideal}；\noteref{cyclic}。原题：《代数学引论》第三版，书页132，扫描PDF第142页。

% 原题逐页视觉核对；扫描PDF第142页；本题证明为本册解答。

\begin{proof}

设$J\lhd R$。对$A=\left(\begin{smallmatrix}a&0\\b&c\end{smallmatrix}\right)\in J$，有
\[
E_{11}AE_{11}=aE_{11},\quad E_{22}AE_{11}=bE_{21},\quad E_{22}AE_{22}=cE_{22},
\]
这些矩阵都在$J$中。分别令$A_1,B_1,C_1$为满足$aE_{11}\in J,bE_{21}\in J,cE_{22}\in J$的整数集合。它们都是$\Z$的加法子群，故分别唯一地形如$m\Z,k\Z,n\Z$，其中$m,k,n\ge0$。

分离公式和加法封闭性共同说明
\[
J=I_{m,k,n}:=\left\{\begin{pmatrix}a&0\\b&c\end{pmatrix}:a\in m\Z,\ b\in k\Z,\ c\in n\Z\right\}.
\]
又因$E_{21}(aE_{11})=aE_{21}$、$(cE_{22})E_{21}=cE_{21}$，必有
\[
m\Z\subseteq k\Z,\qquad n\Z\subseteq k\Z.
\]

反过来，取任意满足上述两包含关系的$m,k,n$。$I_{m,k,n}$对减法封闭。对$X=\left(\begin{smallmatrix}r&0\\s&t\end{smallmatrix}\right)\in R$和$A\in I_{m,k,n}$，
\[
XA=\begin{pmatrix}ra&0\\sa+tb&tc\end{pmatrix},\qquad
AX=\begin{pmatrix}ar&0\\br+cs&ct\end{pmatrix}.
\]
两个对角分量分别仍在$m\Z,n\Z$中；左下分量中$a,c\in k\Z$，故所有项均在$k\Z$中。因此$XA,AX\in I_{m,k,n}$，它确为双边理想。

所以全部理想恰是上述$I_{m,k,n}$。参数条件也可写为$k\mid m$、$k\mid n$，其中$0\mid m$仅当$m=0$；用集合包含写法可统一涵盖零参数。

\end{proof}

\paragraph{复习提示}

先取出各分量，才能排除不同位置的系数之间还存在额外耦合关系。



\subsection{习题17：剩余类矩阵环的单位群阶数}

\label{ex:3:17}

\paragraph{题目}

试计算环$R=\Z/(n)$上$2\times2$全矩阵环的单位群的阶。

\paragraph{考点与思路}

保留中国剩余分解、模$p$约化、两列计数的路线，补齐可逆性提升的理由。

\paragraph{所用结论与原文依据}

\noteref{c3crt}；\noteref{unit}；\exref{1}{44}。原题：《代数学引论》第三版，书页132，扫描PDF第142页。

% 原题逐页视觉核对；扫描PDF第142页；本题证明为本册解答。

\begin{proof}

设$n\ge2$，写$n=\prod_{i=1}^r p_i^{e_i}$。中国剩余定理按矩阵分量给出环同构
\[
M_2(\Z/(n))\cong\bigoplus\limits_{i=1}^r M_2(\Z/(p_i^{e_i})).
\]
分量逐一可逆当且仅当整个元组可逆，所以其单位群是各$\GL_2(\Z/(p_i^{e_i}))$的直积，阶数相乘。

固定$p^e$。交换环上的二阶矩阵$A$可逆当且仅当$\det A$为单位：必要性由$\det A\det A^{-1}=1$得到；充分性由$A\operatorname{adj}A=\operatorname{adj}A\,A=(\det A)E$得到。整数剩余类$a\bmod p^e$为单位当且仅当$p\nmid a$，因为互素时贝祖等式给出逆；被$p$整除时不能与任何整数相乘得到$1$模$p^e$。

因此$A$在模$p^e$下可逆，当且仅当其模$p$像可逆。每个模$p$的矩阵有$p^{4(e-1)}$个提升，四个分量均可独立选择；一个可逆模$p$矩阵的所有这些提升都可逆。于是
\[
|\GL_2(\Z/(p^e))|=p^{4(e-1)}|\GL_2(\F_p)|.
\]
在$\F_p^2$中，可逆矩阵的两列须为有序基。第一列可选任意非零向量，共$p^2-1$种；它张成的直线有$p$个向量，第二列须避开此直线，共$p^2-p$种。因此
\[
\boxed{|\GL_2(\Z/(n))|=\prod_{p^e\parallel n}p^{4(e-1)}(p^2-1)(p^2-p).}
\]
若$n=1$，环只有一个元素，且$0=1$；按单位的双侧逆定义，单位群为单元素群，阶为$1$，也与空乘积约定一致。

\end{proof}

\paragraph{复习提示}

约化后的矩阵可逆，才能保证它的每个提升可逆；理由是行列式为单位的判据。



\subsection{习题18：零因子非空且有限时环必有限}

\label{ex:3:18}

\paragraph{题目}

设$R$为交换环（不必有$1$）。若$R$有零因子但又只有有限个零因子，证明$R$是一个有限环。

\paragraph{考点与思路}

固定一对非零零乘积，以乘法映射把环分成有限多个有限纤维；不预设乘法单位元。

\paragraph{所用结论与原文依据}

\noteref{zerodiv}；\noteref{hom}；\noteref{cosets}。原题：《代数学引论》第三版，书页132，扫描PDF第142页。

% 原题逐页视觉核对；扫描PDF第142页；本题证明为本册解答。

\begin{proof}

设$Z$为$R$中全部零因子的集合。题设说明$Z$非空且有限。因$R$交换，可取$0\ne a,b\in R$使$ab=0$。

考虑加法群同态
\[
L_a:R^+\longrightarrow R^+,\qquad r\longmapsto ar.
\]
其像$aR$有限：对任意$r\in R$，交换性、结合律给出$b(ar)=(ba)r=0$；若$ar\ne0$，则非零元素$b$把$ar$乘为零，所以$ar$为零因子。因此
\[
aR\subseteq Z\cup\{0\}.
\]
其核$K=\{r\in R:ar=0\}$也有限：若$r\in K\setminus\{0\}$，因$a\ne0$，$r$为零因子。所以$K\subseteq Z\cup\{0\}$。

任给$y\in aR$，选$r_y$使$ar_y=y$。对$r\in R$，有
\[
ar=y\ \Longleftrightarrow\ a(r-r_y)=0\ \Longleftrightarrow\ r\in r_y+K.
\]
因此每个非空纤维都是$K$的一个加法陪集，含$|K|$个元素；像只有$|aR|$个元素。于是
\[
|R|=|aR|\,|K|\le(|Z|+1)^2<\infty.
\]
所以$R$是有限环。全部论证只用加法、乘法和交换性，没有使用$1$。

\end{proof}

\paragraph{复习提示}

“有零因子”保证能选到非零$a,b$；无限整环没有零因子，不能据零因子数量有限就判为有限环。



\subsection{习题19：理想关系不具传递性}

\label{ex:3:19}

\paragraph{题目}

设$N$为环$R$（不必有$1$）的理想，而且$I$又为$N$（作为一个环）的理想，问$I$是不是$R$的理想？

\paragraph{考点与思路}

沿用严格上三角矩阵的反例，写明每个乘积落在哪里及哪一个乘积破坏吸收性。

\paragraph{所用结论与原文依据}

\noteref{ideal}。原题：《代数学引论》第三版，书页132，扫描PDF第142页。

% 原题逐页视觉核对；扫描PDF第142页；本题证明为本册解答。

\begin{proof}

设$F$为域，取
\[
R=\left\{\begin{pmatrix}0&a&b\\0&0&c\\0&0&0\end{pmatrix}:a,b,c\in F\right\},\quad
N=FE_{12}+FE_{13},\quad I=FE_{12}.
\]
对$X=aE_{12}+bE_{13}+cE_{23}\in R$和$Y=uE_{12}+vE_{13}\in N$，矩阵单位乘法$E_{ij}E_{k\ell}=\delta_{jk}E_{i\ell}$给出
\[
XY=0,\qquad YX=ucE_{13}\in N.
\]
加法与取负的封闭性由线性表达式得到，所以$N\lhd R$。同时$N^2=0$，因此$I$对$N$的左右乘法均吸收为零，又是加法子群，故$I\lhd N$。

但是$E_{12}\in I$、$E_{23}\in R$，而
\[
E_{12}E_{23}=E_{13}\notin I.
\]
所以$I$不是$R$的理想。此处$R$没有幺元，符合一般环不预设幺元的约定。

\end{proof}

\paragraph{复习提示}

验证$I\lhd N$只检查$N$中元素，验证$I\lhd R$则必须检查整个$R$中的元素。



\subsection{习题20：非单位属于某个极大理想}

\label{ex:3:20}

\paragraph{题目}

设$R$为一个交换环，$a$是$R$的不可逆元。证明$R$有一个极大理想$M$而且包含$a$。

\paragraph{考点与思路}

从真主理想$(a)$出发，以佐恩引理选极大元，并检查它在全部真理想中也极大。

\paragraph{所用结论与原文依据}

\noteref{c3prime}；\noteref{c3principal}；\noteref{c3maxexists}。原题：《代数学引论》第三版，书页132，扫描PDF第142页。

% 原题逐页视觉核对；扫描PDF第142页；本题证明为本册解答。

\begin{proof}

这里$R$为交换幺环，$a$为非单位。若$(a)=R$，则$1=ra$对某个$r\in R$成立；交换性使$r$为$a$的双侧逆，矛盾。因此$(a)\ne R$。

令$\mathcal S=\{I\lhd R:(a)\subseteq I\ne R\}$，以包含关系赋序。它含$(a)$。设$\mathcal C\subseteq\mathcal S$为非空链，令$U=\bigcup_{I\in\mathcal C}I$。任取$x,y\in U$，可在链中找一个同时含二者的理想，故$x-y\in U$；对$r\in R$，$rx\in U$。所以$U$是包含$(a)$的理想。

若$U=R$，则$1\in U$，从而链中某个$I$含$1$，使$I=R$，矛盾。因此$U\in\mathcal S$，是该链的上界。空链可取$(a)$为上界。由佐恩引理，$\mathcal S$有极大元$M$。

若$M\subsetneq J\subsetneq R$且$J$为理想，则$J$也含$(a)$，故$J\in\mathcal S$，违背$M$的极大性。因此$M$为$R$的极大理想，并含$a$。

\end{proof}

\paragraph{复习提示}

链的并仍是真理想，关键在于一旦并含$1$，某个链成员就已含$1$。



\subsection{习题21：雅各布森根元素的可逆性性质}

\label{ex:3:21}

\paragraph{题目}

设$R$为一交换环，$J(R)$表示$R$的全部极大理想的交，$J(R)$叫做环$R$的雅各布森（Jacobson）根。证明$J(R)$的每个元素$a$有如下性质：对所有$x\in R$，$1-ax$都是可逆元。（与第三章定理9推论2、第一章习题38作比较。）

\paragraph{考点与思路}

若$1-ax$不是单位，先把它放入极大理想，再与$ax$相加得到矛盾。

\paragraph{所用结论与原文依据}

\noteref{c3prime}；\noteref{c3nil}；\exref{3}{20}；\exref{1}{38}。原题：《代数学引论》第三版，书页133，扫描PDF第143页。

% 原题逐页视觉核对；扫描PDF第143页；本题证明为本册解答。

\begin{proof}

任取$a\in J(R)$和$x\in R$。反设$1-ax$不可逆。由\exref{3}{20}，存在极大理想$M$包含$1-ax$。由于$J(R)$是全部极大理想的交，$a\in M$；理想的吸收性给出$ax\in M$。于是
\[
1=(1-ax)+ax\in M.
\]
含$1$的理想等于整个环，与$M$是真理想矛盾。因此$1-ax$可逆。因$x$任意，题述性质对所有$x\in R$成立。

与原书所指结论比较：第三章定理9推论2把诣零根表示为全部素理想的交，而这里取的是全部极大理想的交；第一章习题38从$a$幂零出发给出有限等比和逆元。此题把“对所有$x$，$1-ax$可逆”从幂零元推广到$J(R)$中的元素，但没有预先断言这些元素都幂零。

\end{proof}

\paragraph{复习提示}

本题先证$J(R)$中元素具有该性质；反方向正是下一题。



\subsection{习题22：强拟正则元属于雅各布森根}

\label{ex:3:22}

\paragraph{题目}

设$R$为一交换环。如果元素$a\in R$具有习题21中的性质，则$a$叫做一个强拟正则元。证明$R$的强拟正则元都属于$J(R)$。（参看习题21。）

\paragraph{考点与思路}

对一个不包含$a$的极大理想，在其商域中选择$a$的逆，并提升到环内。

\paragraph{所用结论与原文依据}

\noteref{c3prime}；\exref{3}{21}。原题：《代数学引论》第三版，书页133，扫描PDF第143页。

% 原题逐页视觉核对；扫描PDF第143页；本题证明为本册解答。

\begin{proof}

设$a$是强拟正则元，即对每个$x\in R$，$1-ax$均可逆。反设$a\notin J(R)$。按交集定义，存在极大理想$M$使$a\notin M$。商环$R/M$为域，非零类$a+M$有逆。选取逆元的代表$x\in R$，则
\[
(a+M)(x+M)=1+M,
\qquad 1-ax\in M.
\]
但$1-ax$是单位，若它在$M$中，乘以其逆便得$1\in M$，与$M\ne R$矛盾。因此$a\in J(R)$，所以全部强拟正则元都属于$J(R)$。

\end{proof}

\paragraph{复习提示}

逆元是在商域中取的，不能直接假设$a$在原环中已经可逆。



\subsection{习题23：有限交换环的素理想必极大}

\label{ex:3:23}

\paragraph{题目}

设$R$为一交换环。证明：若$R$有限，则$R$的素理想都是极大理想。

\paragraph{考点与思路}

取素理想的商整环，用有限集的消去律证明非零元素都有逆。

\paragraph{所用结论与原文依据}

\noteref{c3prime}；\noteref{zerodiv}；\noteref{field}。原题：《代数学引论》第三版，书页133，扫描PDF第143页。

% 原题逐页视觉核对；扫描PDF第143页；本题证明为本册解答。

\begin{proof}

设$P$为有限交换环$R$的素理想。由素理想的商环判据，$D=R/P$为整环；特别地，$D$非零，且是有限集合。

任取$0\ne a\in D$。若$ax=ay$，无零因子性给出$x=y$，故$L_a:x\mapsto ax$为有限集合$D$上的单射，因而满射。于是存在$b\in D$使$ab=1_D$。交换性给出$ba=1_D$。因此每个非零元均有逆，$D$为域。

由极大理想的商域判据，$P$为极大理想。

\end{proof}

\paragraph{复习提示}

本题用的是有限整环为域；不能引用原书第18题来代替，它讨论的是有有限多个零因子的环。



\subsection{习题24：环的内直和定理与商环直和定理的等价性}

\label{ex:3:24}

\paragraph{题目}

证明第三章$\S2$的定理4和定理$5'$等价。

\paragraph{考点与思路}

按原书§2书页107的说明，正向取其余理想的交，反向取其余直和分量的和，分别核对定理的条件。

\paragraph{所用结论与原文依据}

\noteref{c3internal}；\noteref{c3direct}；\noteref{c3comaximal}；\noteref{c3isomorphism}。原题：《代数学引论》第三版，书页133，扫描PDF第143页。

% 原题逐页视觉核对；扫描PDF第143页；本题证明为本册解答。

\begin{proof}

原书第三章§2的定理4（书页104）是环的内直和判别：若$R_i\lhd R$、$R=\sum_iR_i$，且
\[
R_i\cap\sum_{j\ne i}R_j=0\qquad(1\le i\le r),
\]
则$R\cong\bigoplus\limits_{i=1}^rR_i$。定理$5'$（书页107）是：若$N_i\lhd R$两两互素且$\bigcap\limits_{i=1}^rN_i=0$，则$R\cong\bigoplus\limits_{i=1}^rR/N_i$。以下证明二者互相推出，不使用待证的另一条作为前提。

\begin{enumerate}
\item 由定理4推出定理$5'$。设$N_i$两两互素，且交为零，令
\[
R_i=\bigcap_{j\ne i}N_j.
\]
当$r=1$时取$R_1=R$，结论即$R\cong R/0$。以下设$r\ge2$。理想的交仍是理想，故各$R_i\lhd R$。

先证$R_i+N_i=R$。令$M_i$为除$N_i$外其余$N_j$按下标顺序的理想积。由原书§2的互素理想乘积引理反复应用，$M_i+N_i=R$；又$M_i\subseteq R_i$，故$R_i+N_i=R$。可取$e_i\in R_i$，使$1-e_i\in N_i$。

对固定$k$，$i\ne k$时$e_i\in R_i\subseteq N_k$，因而$1-\sum_i e_i\in N_k$。此对每个$k$都成立，交为零使$1=\sum_i e_i$。任意$x\in R$有$x=\sum_i xe_i$，且$xe_i\in R_i$，故$R=\sum_iR_i$。

又$\sum_{j\ne i}R_j\subseteq N_i$，所以
\[
R_i\cap\sum_{j\ne i}R_j\subseteq R_i\cap N_i=\bigcap_kN_k=0.
\]
定理4的全部条件成立，得到$R\cong\bigoplus_iR_i$。自然同态在$R_i$上的限制$R_i\to R/N_i$，核为$R_i\cap N_i=0$，且由$R_i+N_i=R$知它满射，故$R_i\cong R/N_i$。合并这些同构，即得定理$5'$。

\item 由定理$5'$推出定理4。设$R_i$满足定理4的条件，令
\[
N_i=\sum_{j\ne i}R_j.
\]
每个$N_i$为理想。对$i\ne j$，$N_i+N_j$包含全部$R_k$，故等于$R$，所以这些理想两两互素。

题设中的交为零条件先给出加法分解唯一性：若$\sum_i a_i=\sum_i b_i$且$a_i,b_i\in R_i$，则$a_i-b_i=-\sum_{j\ne i}(a_j-b_j)\in R_i\cap N_i=0$，故$a_i=b_i$。若$x\in\bigcap_iN_i$，在其唯一分解$x=\sum_i x_i$中，$x\in N_i$使第$i$分量为零；所有分量都为零，故$x=0$。

由定理$5'$，$R\cong\bigoplus_iR/N_i$。对每个$i$，自然同态$R_i\to R/N_i$的核为$R_i\cap N_i=0$；又$R=R_i+N_i$，故满射，得到$R_i\cong R/N_i$。因此$R\cong\bigoplus_iR_i$，即定理4。
\end{enumerate}

\end{proof}

\paragraph{复习提示}

题目中的“第三章§2”是环的直和；定理$5'$的撇号必须保留。



\subsection{习题25：单位根群的满而非单自同态}

\label{ex:3:25}

\paragraph{题目}

设$G$是复数域中全部单位根组成的乘法群。请具体写出$G$的一个自同态，它是满的但不是单一的。

\paragraph{考点与思路}

保留平方映射，取平方根证明满射，用$1$和$-1$证明不单。

\paragraph{所用结论与原文依据}

\noteref{hom}。原题：《代数学引论》第三版，书页133，扫描PDF第143页。

% 原题逐页视觉核对；扫描PDF第143页；本题证明为本册解答。

\begin{proof}

记$G=\{z\in\C:z^m=1\text{ 对某个 }m\ge1\}$，定义$\varphi:G\to G$为$\varphi(z)=z^2$。
若$z^m=1$，则$(z^2)^m=1$，所以$\varphi$确实映入$G$。因复数乘法交换，
\[
\varphi(zw)=(zw)^2=z^2w^2=\varphi(z)\varphi(w),
\]
它是自同态。

任取$w\in G$，取$m\ge1$使$w^m=1$。复数中存在$z$使$z^2=w$，于是$z^{2m}=w^m=1$，说明$z\in G$，且$\varphi(z)=w$。因此$\varphi$满射。
但$1\ne-1$均在$G$中，且$\varphi(1)=\varphi(-1)=1$，所以$\varphi$不单射。

\end{proof}

\paragraph{复习提示}

取出的平方根还须属于单位根群，而不只是属于复数域。



以下26题至29题可参考华罗庚、万哲先著的《典型群》一书第一章§4。

\subsection{习题26：若尔当同态的三项性质}

\label{ex:3:26}

\paragraph{题目}

如果环$R$到环$R'$的映射$\eta$满足
\begin{enumerate}[label=(\roman*)]
\item $\eta$是加法群同态，即$\eta(a+b)=\eta(a)+\eta(b)$；
\item $\eta(1)=1'$；
\item $\eta(aba)=\eta(a)\eta(b)\eta(a)$，
\end{enumerate}
则$\eta$叫做一个若尔当同态。证明：若$\eta$是一个若尔当同态，则$\eta$具有如下性质：
\begin{enumerate}[label=(\roman*)]
\item $\eta(a^k)=\eta(a)^k$；
\item $\eta(abc+cba)=\eta(a)\eta(b)\eta(c)+\eta(c)\eta(b)\eta(a)$；
\item $\eta(ab+ba)=\eta(a)\eta(b)+\eta(b)\eta(a)$。
\end{enumerate}
显然环的同态和反同态都是若尔当同态。

\paragraph{考点与思路}

按原书定义先验证全部非负整数次幂，再对三重乘积展开，并取中间变量为1得到最后一项。

\paragraph{所用结论与原文依据}

\noteref{ringlaws}；\noteref{unit}。原题：《代数学引论》第三版，书页133，扫描PDF第143页。

% 原题逐页视觉核对；扫描PDF第143页；本题证明为本册解答。

\begin{proof}

\begin{enumerate}[label=(\roman*)]
\item 对非负整数$k$证明。由保幺性，$\eta(a^0)=1'=\eta(a)^0$；$k=1$为恒等式。若已知$\eta(a^m)=\eta(a)^m$，定义中的三重乘积性质给出
\[
\eta(a^{m+2})=\eta(aa^ma)=\eta(a)\eta(a^m)\eta(a)=\eta(a)^{m+2}.
\]
分别从$m=0,1$开始归纳，得到全部偶数次幂与奇数次幂，故$\eta(a^k)=\eta(a)^k$对$k\ge0$成立。若另假设$a$为单位，取$b=a^{-2}$在$aba=1$中，可得$\eta(a)\eta(a^{-2})\eta(a)=1'$，从而$\eta(a)$既有左逆又有右逆，二者相等；再把$aa^{-1}a=a$取像并消去$\eta(a)$，得$\eta(a^{-1})=\eta(a)^{-1}$。所以对已有定义的负整数次幂，公式也成立。
\item 将$a+c$代入若尔当同态定义的两个外侧位置，得
\[
\eta((a+c)b(a+c))=(\eta(a)+\eta(c))\eta(b)(\eta(a)+\eta(c)).
\]
展开左侧乘积并使用加法性，右侧用分配律展开，然后分别消去$\eta(aba)=\eta(a)\eta(b)\eta(a)$及$\eta(cbc)=\eta(c)\eta(b)\eta(c)$，得到
\[
\eta(abc+cba)=\eta(a)\eta(b)\eta(c)+\eta(c)\eta(b)\eta(a).
\]
\item 在第二项中将中间变量取为$1$，外侧变量取为$a,b$，由$\eta(1)=1'$，得到
\[
\eta(ab+ba)=\eta(a)\eta(b)+\eta(b)\eta(a).
\]
\end{enumerate}
最后验证原题末句：保幺同态对$aba$取像显然得到所需三重乘积；保幺反同态则给出$\eta(aba)=\eta(a)\eta(b)\eta(a)$，因为反序后外侧仍同为$a$。结合加法性和保幺性，两者均是题中定义的若尔当同态。

\end{proof}

\paragraph{复习提示}

三重乘积公式是定义的一部分；证明不除以2，适用于特征2。



\subsection{习题27：若尔当同态对一对元素的二择一性质}

\label{ex:3:27}

\paragraph{题目}

（雅各布森--瑞卡特（Jacobson--Rickart））若$\eta$是环$R$到无零因子的环$R'$的一个若尔当同态，证明对任意$a,b\in R$，$\eta(ab)=\eta(a)\eta(b)$和$\eta(ab)=\eta(b)\eta(a)$有一个成立。

\paragraph{考点与思路}

对三重乘积极化式代入$ab,b,a$，再按固定次序因式分解。

\paragraph{所用结论与原文依据}

\noteref{zerodiv}；\exref{3}{26}。原题：《代数学引论》第三版，书页133，扫描PDF第143页。

% 原题逐页视觉核对；扫描PDF第143页；本题证明为本册解答。

\begin{proof}

沿用\exref{3}{26}的若尔当同态约定，记
\[
A=\eta(a),\quad B=\eta(b),\quad C=\eta(ab).
\]
在该题的极化式中令外侧变量为$ab,a$，中间变量为$b$，得到
\[
\eta(ab^2a+abab)=CBA+ABC.
\]
由三重乘积公式及平方公式，
\[
\eta(ab^2a)=A\eta(b^2)A=AB^2A,\qquad
\eta(abab)=\eta((ab)^2)=C^2.
\]
加法性给出$C^2+AB^2A=CBA+ABC$。保持因子顺序展开，有
\begin{align*}
(C-AB)(C-BA)
&=C^2-CBA-ABC+ABBA\\
&=C^2-CBA-ABC+AB^2A=0.
\end{align*}
目标环无零因子，故$C-AB=0$或$C-BA=0$。即对这对$a,b$，
\[
\eta(ab)=\eta(a)\eta(b)\quad\text{或}\quad\eta(ab)=\eta(b)\eta(a).
\]

\end{proof}

\paragraph{复习提示}

本题的二择一暂时依赖于$a,b$；还未证明所有元素对都选择同一乘法方向。



\subsection{习题28：逐对二择一推出全局同态或反同态}

\label{ex:3:28}

\paragraph{题目}

（华罗庚定理）设$\eta$为环$R$到环$R'$的一个映射，满足
\begin{enumerate}[label=(\roman*)]
\item $\eta(a+b)=\eta(a)+\eta(b)$；
\item $\eta(1)=1'$；
\item $\eta(ab)=\eta(a)\eta(b)$或$\eta(ab)=\eta(b)\eta(a)$，
\end{enumerate}
则$\eta$是一个同态或反同态。（后者定义为$R$到$R'$的映射$\eta$满足(i)、(ii)和对于所有$a,b\in R$，总有$\eta(ab)=\eta(b)\eta(a)$。）

\paragraph{考点与思路}

对第一个变量、第二个变量分别使用一次“两子群之并”论证。

\paragraph{所用结论与原文依据}

\noteref{subgroup}；\noteref{c3anti}。原题：《代数学引论》第三版，书页133，扫描PDF第143页。

% 原题逐页视觉核对；扫描PDF第143页；本题证明为本册解答。

\begin{proof}

先证所用群论事实：若群$G=H\cup K$且$H,K$均为子群，则$H=G$或$K=G$。否则二者都真，且互不包含，可取$h\in H\setminus K$、$k\in K\setminus H$。若$hk\in H$则$k=h^{-1}(hk)\in H$；若$hk\in K$则$h=(hk)k^{-1}\in K$，两种都矛盾。因此$hk$不属于并集，违背$G=H\cup K$。

固定$b\in R$，定义
\begin{align*}
H_b&=\{a:\eta(ab)=\eta(a)\eta(b)\},\\
K_b&=\{a:\eta(ab)=\eta(b)\eta(a)\}.
\end{align*}
二者都含$0$。若$a,c\in H_b$，加法性与分配律给出
\[
\eta((a-c)b)=\eta(ab)-\eta(cb)
=(\eta(a)-\eta(c))\eta(b)=\eta(a-c)\eta(b),
\]
所以$a-c\in H_b$。若$a,c\in K_b$，同样计算为
\[
\eta((a-c)b)=\eta(b)(\eta(a)-\eta(c))=\eta(b)\eta(a-c),
\]
故$K_b$也为加法子群。题设说明$R=H_b\cup K_b$，因而每个$b$都满足$H_b=R$或$K_b=R$。

令$B_+=\{b:H_b=R\}$、$B_-=\{b:K_b=R\}$。二者含$0$。若$b,c\in B_+$，则对所有$a$，
\[
\eta(a(b-c))=\eta(ab)-\eta(ac)=\eta(a)(\eta(b)-\eta(c))=\eta(a)\eta(b-c),
\]
所以$b-c\in B_+$。若$b,c\in B_-$，则
\[
\eta(a(b-c))=\eta(b)\eta(a)-\eta(c)\eta(a)=\eta(b-c)\eta(a),
\]
故$b-c\in B_-$。于是$B_+,B_-$都是加法子群，且$R=B_+\cup B_-$。再次使用所证事实，$B_+=R$或$B_-=R$。

第一种情形给出对所有$a,b$都有$\eta(ab)=\eta(a)\eta(b)$；第二种给出对所有$a,b$都有$\eta(ab)=\eta(b)\eta(a)$。结合既有加法性，$\eta$分别为同态或反同态。

\end{proof}

\paragraph{复习提示}

使用两次子群并，才能消除方向分别依赖于两个变量的可能。



\subsection{习题29：环的若尔当同态是同态或反同态}

\label{ex:3:29}

\paragraph{题目}

（雅各布森--瑞卡特）证明环$R$到无零因子的环$R'$的任一个若尔当同态必是一个同态或反同态。

\paragraph{考点与思路}

把第27题的逐对结论送入第28题；原题源环是一般环，不要求为除环。

\paragraph{所用结论与原文依据}

\exref{3}{26}；\exref{3}{27}；\exref{3}{28}。原题：《代数学引论》第三版，书页133，扫描PDF第143页。

% 原题逐页视觉核对；扫描PDF第143页；本题证明为本册解答。

\begin{proof}

若$\eta:R\to R'$为若尔当同态，依\exref{3}{26}的定义，它是加法群同态，且$\eta(1)=1'$、$\eta(aba)=\eta(a)\eta(b)\eta(a)$。

因为目标环$R'$无零因子，\exref{3}{27}给出：对任意$a,b\in R$，$\eta(ab)$等于$\eta(a)\eta(b)$或$\eta(b)\eta(a)$。

因此$\eta$满足\exref{3}{28}华罗庚定理的全部条件，由该题结论，$\eta$是同态或反同态。这里从始至终没有假设$R$为除环，也没有把仅保平方代替原书的若尔当同态定义。

\end{proof}

\paragraph{复习提示}

“加法、保幺、保三重乘积”就是原书第26题的定义，不需要另设一套定义。



\subsection{习题30：模素数幂环的单位、幂零元与唯一素理想}

\label{ex:3:30}

\paragraph{题目}

设$p$为一素数，$n$为任一大于$1$的整数，$R=\Z/(p^n)$。证明：
\begin{enumerate}[label=(\roman*)]
\item $R$的元素不是单位便是零幂元；
\item $R$恰有一个素理想，记作$P$；
\item 商环$R/P$是一个域。
\end{enumerate}

\paragraph{考点与思路}

按原题三项顺序，先用互素判别单位与幂零，再列出理想确定唯一素理想，最后求商域。

\paragraph{所用结论与原文依据}

\noteref{c3correspondence}；\noteref{c3prime}；\noteref{unit}；\noteref{c3nil}。原题：《代数学引论》第三版，书页133，扫描PDF第143页。

% 原题逐页视觉核对；扫描PDF第143页；本题证明为本册解答。

\begin{proof}

\begin{enumerate}[label=(\roman*)]
\item 将$a+(p^n)$记为$\bar a$。若$p\nmid a$，贝祖等式给出整数$s,t$使$sa+tp^n=1$，所以$\bar s\bar a=1$，$\bar a$为单位。

若$p\mid a$，则$p^n\mid a^n$，故$\bar a^{\,n}=0$，$\bar a$为幂零元。这包括$\bar a=0$。因$n>1$，环非零；单位不可能幂零，否则乘以其逆的相应幂会得到$1=0$。因此两类穷尽全部元素。
\item 由$\Z\to R$的理想对应，$R$的理想恰为$P_k=(p^k)/(p^n)$，其中$0\le k\le n$。$P_0=R$不是真理想，因而不素。对$2\le k\le n$，$\bar p\notin P_k$、$\overline{p^{k-1}}\notin P_k$，但二者乘积在$P_k$中，所以$P_k$不素。

令$P=P_1$。若$\bar a\bar b\in P$，则$p\mid ab$，因$p$为素数，$p\mid a$或$p\mid b$，即$\bar a\in P$或$\bar b\in P$。又$P\ne R$，故$P$为素理想。由上面全部候选的排除，它是唯一素理想。
\item 模$p$的约化映射
\[
\pi:R\longrightarrow\Z/(p),\qquad a+(p^n)\longmapsto a+(p)
\]
良定义且为满环同态，其核恰为$P$。由同态基本定理，$R/P\cong\Z/(p)$，后者为域。
\end{enumerate}

\end{proof}

\paragraph{复习提示}

第一项原题说的是“单位或幂零元”，不是“单位或零因子”；零元也已被包含在幂零元一类。



\subsection{习题31：模素数幂系数多项式的单位、零因子与幂零元}

\label{ex:3:31}

\paragraph{题目}

设$p$为一素数，$n$为一整数，$R=\Z/(p^n)$，试具体指出多项式环$R[x]$中哪些元素是单位、零因子或是零幂元。

\paragraph{考点与思路}

保留模$p$与系数幂零的判别；为避免跨题缺口，在题内证明单位的必要条件，并展开最低赋值系数论证。

\paragraph{所用结论与原文依据}

\noteref{c3polyquot}；\noteref{c3degree}；\exref{3}{30}。原题：《代数学引论》第三版，书页133，扫描PDF第143页。

% 原题逐页视觉核对；扫描PDF第143页；本题证明为本册解答。

\begin{proof}

设$n\ge1$、$f=\sum_{i=0}^m a_ix^i\in(\Z/(p^n))[x]$。先求单位。若$fg=1$，比较常数项知$a_0$为单位。逐系数模$p$约化后，在整环$\F_p[x]$中有$\bar f\bar g=1$。次数可加，所以$\bar f$为非零常数；因此$p\nmid a_0$，而所有$i>0$的$a_i$均被$p$整除。

反过来，若满足这些系数条件，$a_0$可逆，令$h=a_0^{-1}(f-a_0)$。$h$的每个系数被$p$整除，所以$h^n=0$。于是
\[
f^{-1}=a_0^{-1}\sum_{j=0}^{n-1}(-h)^j,
\]
左右相乘均由有限等比和得到$1$。故
\[
f\text{ 可逆}\ \Longleftrightarrow\ p\nmid a_0\text{ 且 }p\mid a_i\ (i>0).
\]

再求非零零因子。若$f\ne0$且全部系数均被$p$整除，则常数$\overline{p^{n-1}}\ne0$且$\overline{p^{n-1}}f=0$，故$f$为零因子（$n=1$时不存在这样的非零$f$）。

反过来，设$f$有不被$p$整除的系数，取最小的这类下标$j$。对剩余类$r$，令$v(r)$为最大的$s\in\{0,\ldots,n\}$，使$r$被$p^s$整除，其中$v(0)=n$。任取非零$g=\sum b_kx^k$，令$s=\min_kv(b_k)<n$，再取达到该最小值的最小下标$k$。

在$fg$的$x^{j+k}$项中，$a_jb_k$的赋值恰为$s$。其余乘积$a_ib_\ell$满足$i+\ell=j+k$。若$i<j$，则$v(a_i)\ge1$而$v(b_\ell)\ge s$；若$i>j$，则$\ell<k$，由$k$的选择有$v(b_\ell)\ge s+1$。故所有其他项都被$p^{s+1}$整除，$a_jb_k$却不被$p^{s+1}$整除。总和不能为零模$p^n$，因$s<n$。所以$fg\ne0$。

因此非零零因子恰为全部系数被$p$整除的非零多项式。


最后求幂零元。若$f$的全部系数均被$p$整除，则$f^n=0$，因为展开后的每项都含$n$个$p$的倍数。反过来，若$f^N=0$，逐系数模$p$约化后，$\bar f^N=0$在整环$\F_p[x]$中成立，故$\bar f=0$，即每个系数都被$p$整除。因此
\[
f\text{ 幂零}\quad\Longleftrightarrow\quad p\mid a_i\text{ 对全部 }i.
\]
所以零因子为上述幂零多项式中的非零元素，而幂零元一类还包括零多项式。这里按$p^n$为正素数幂即$n\ge1$理解；若取$n=0$，系数环为零环，唯一多项式为零，它是幂零元并按$0=1$的约定为单位，没有非零零因子。

\end{proof}

\paragraph{复习提示}

单位、零因子、幂零元是原题要求的三类；后两类在此仅相差是否包含零多项式。



\subsection{习题32：直和环的理想分解及整数商环的根与单位群}

\label{ex:3:32}

\paragraph{题目}

设$R=R_1\oplus R_2\oplus\cdots\oplus R_r$是环$R$的一个内直和，$N$为$R$的一个理想。证明
\[
N=(N\cap R_1)\oplus(N\cap R_2)\oplus\cdots\oplus(N\cap R_r),
\]
其中$N\cap R_i$有可能等于$(0)$。设$n>1$，试决定$R=\Z/(n)$的最大零幂理想和单位群。

\paragraph{考点与思路}

用各分量的单位元投影理想，再按中国剩余分解逐个研究模素数幂因子。

\paragraph{所用结论与原文依据}

\noteref{c3direct}；\noteref{c3crt}；\noteref{c3nil}；\exref{3}{30}。原题：《代数学引论》第三版，书页133--134，扫描PDF第143页起。

% 原题逐页视觉核对；扫描PDF第143页；本题证明为本册解答。

\begin{proof}

\begin{enumerate}
\item 设每个$R_i$为幺环，并将它等同于$R=\bigoplus\limits_{i=1}^rR_i$中仅第$i$分量非零的理想。令$e_i$为该分量的单位元，其他分量为零。对$x=(x_1,\ldots,x_r)\in N$，理想吸收性给出
\[
e_ix=(0,\ldots,x_i,\ldots,0)\in N\cap R_i,\qquad
x=\sum_{i=1}^r e_ix.
\]
所以$N\subseteq\sum_i(N\cap R_i)$；反向包含来自每个交均含于$N$。不同坐标的表示唯一，因此
\[
N=\bigoplus\limits_{i=1}^r(N\cap R_i).
\]
若不要求各分量有幺元，取两个零乘法环$\Z$的直和，斜子群$\{(m,m):m\in\Z\}$就是反例。
\item 设$n\ge2$，写$n=\prod_i p_i^{a_i}$，令$d=\prod_i p_i$。中国剩余定理给出$\Z/(n)\cong\bigoplus_i\Z/(p_i^{a_i})$。在第$i$个因子中，$p_i$的倍数取$a_i$次幂即为零；若不被$p_i$整除，则是单位，非零环中的单位不可能幂零。因此幂零元恰为每个分量中被$p_i$整除的元。整数$a$的剩余类满足此条件当且仅当$d\mid a$。

故幂零元理想为$N=(d)/(n)$。它本身是幂零理想：取$A=\max_i a_i$，因$n\mid d^A$，有$N^A=0$。任何幂零理想的每个元素都幂零，故包含于$N$。因此$N$确为最大的幂零理想。
\item 分量乘积为$1$等价于每个分量乘积为$1$，故
\[
(\Z/(n))^\times\cong\prod_i(\Z/(p_i^{a_i}))^\times.
\]
在模$p^a$环中，不是单位的元素恰为$p$的倍数，共$p^{a-1}$个，所以单位有$p^a-p^{a-1}$个。因此
\[
|(\Z/(n))^\times|=\prod_i(p_i^{a_i}-p_i^{a_i-1})
=n\prod_{p\mid n}\left(1-\frac1p\right).
\]
以上直积分解给出单位群；并未无依据地把每个因子再断言为循环群。$n=1$时零环的最大幂零理想为零理想，单位群为单元素群。
\end{enumerate}

\end{proof}

\paragraph{复习提示}

一般幂零元理想未必幂零；本题额外证明$N^A=0$，才得到最大幂零理想。



\subsection{习题33：三次多项式在六元环中的六个根}

\label{ex:3:33}

\paragraph{题目}

证明$x^3-x$在$\Z/(6)$内有$6$个根。

\paragraph{考点与思路}

沿用连续三个整数的整除性，并说明取代表元不影响代入值。

\paragraph{所用结论与原文依据}

\noteref{c3roots}；\noteref{c3poly}。原题：《代数学引论》第三版，书页134，扫描PDF第144页。

% 原题逐页视觉核对；扫描PDF第144页；本题证明为本册解答。

\begin{proof}

对任意整数$a$，
\[
a^3-a=(a-1)a(a+1).
\]
三个连续整数中至少一个为偶数，且恰有一个被$3$整除，所以$2\mid a^3-a$且$3\mid a^3-a$。因$2,3$互素，$6\mid a^3-a$。

因此对每个$\bar a\in\Z/(6)$，都有$\bar a^3-\bar a=0$；模$6$的代入只依赖于剩余类，因为自然同态保加法和乘法。于是全部六个元素$\bar0,\bar1,\ldots,\bar5$都是根，且环中没有其他元素，恰有六个根。

\end{proof}

\paragraph{复习提示}

次数限制根数的结论要求系数环为整环；$\Z/(6)$有零因子。



\subsection{习题34：四元数中的二次方程有无穷多个根}

\label{ex:3:34}

\paragraph{题目}

证明$x^2+1$在四元数体$H$内有无穷多个根。

\paragraph{考点与思路}

按原习题解展开纯虚四元数的平方，给出明确的无限族来证明互异性。

\paragraph{所用结论与原文依据}

\noteref{quaternion}；\noteref{c3qroots}。原题：《代数学引论》第三版，书页134，扫描PDF第144页。

% 原题逐页视觉核对；扫描PDF第144页；本题证明为本册解答。

\begin{proof}

取纯虚四元数$q=bI+cJ+dK$。由乘法表，$IJ+JI=IK+KI=JK+KJ=0$，所以
\begin{align*}
q^2&=b^2I^2+c^2J^2+d^2K^2
+bc(IJ+JI)+bd(IK+KI)+cd(JK+KJ)\\
&=-(b^2+c^2+d^2).
\end{align*}
因此每个满足$b^2+c^2+d^2=1$的三元组都给出$x^2+1$的根。

例如对每个$t\in[0,1]$，令
\[
q_t=\frac{1-t^2}{1+t^2}I-\frac{2t}{1+t^2}K.
\]
其系数平方和为$((1-t^2)^2+4t^2)/(1+t^2)^2=1$，故$q_t^2=-1$。若$0\le s<t\le1$，则
\[
\frac{1-s^2}{1+s^2}-\frac{1-t^2}{1+t^2}
=\frac{2(t^2-s^2)}{(1+s^2)(1+t^2)}>0,
\]
所以$q_s\ne q_t$。因此有无穷多个根。这个族也等于$(1+tJ)I(1+tJ)^{-1}$，与原笔记的共轭构造一致。

\end{proof}

\paragraph{复习提示}

即使除环没有零因子，乘法不交换时也不能直接套用整环上的根数定理。



\subsection{习题35：两个平方根之和的极小多项式}

\label{ex:3:35}

\paragraph{题目}

在复数域内，$u=\sqrt2+\sqrt3$是有理数域$\Q$上的代数元，求$u$的极小多项式。

\paragraph{考点与思路}

先求四次关系，再从$u$恢复两个平方根，并用线性无关性证明极小次数恰为四。

\paragraph{所用结论与原文依据}

\noteref{c3minimal}；\noteref{c3division}。原题：《代数学引论》第三版，书页134，扫描PDF第144页。

% 原题逐页视觉核对；扫描PDF第144页；本题证明为本册解答。

\begin{proof}

令$u=\sqrt2+\sqrt3$。有
\[
u^2=5+2\sqrt6,\qquad (u^2-5)^2=24,
\]
故$u^4-10u^2+1=0$。候选多项式为$f=x^4-10x^2+1$。

先验证$\sqrt3\notin\Q(\sqrt2)$。若$\sqrt3=r+s\sqrt2$，$r,s\in\Q$，平方并利用$\sqrt2$无理，得
\[
3=r^2+2s^2,
\qquad 2rs=0.
\]
若$s=0$，则$r^2=3$；若$r=0$，则$s^2=3/2$。二者均不可能为有理数的平方，因为有理数平方的任一素因子指数均为偶数，而$3$在这两个数中的指数均为奇数。于是$1,\sqrt2,\sqrt3,\sqrt6$在$\Q$上线性无关：若
\[
(a+b\sqrt2)+(c+d\sqrt2)\sqrt3=0,
\]
且$c+d\sqrt2\ne0$，可在域$\Q(\sqrt2)$中相除得到$\sqrt3\in\Q(\sqrt2)$，矛盾；故$c=d=0$，继而$a=b=0$。

又$(\sqrt3+\sqrt2)(\sqrt3-\sqrt2)=1$，所以$u^{-1}=\sqrt3-\sqrt2$，从而
\[
\sqrt3=\frac{u+u^{-1}}2,\qquad
\sqrt2=\frac{u-u^{-1}}2.
\]
由$f(u)=0$还可得$u^{-1}=10u-u^3\in\Q[u]$，故$\Q[u]$含上述四个线性无关元素。另一方面，$u^4=10u^2-1$使每个$u$的高次幂都化为$1,u,u^2,u^3$的有理线性组合。

若存在次数$d<4$的非零多项式$g$使$g(u)=0$，用其非零首项系数相除后，可以将$u^d$及全部更高幂降为$1,u,\ldots,u^{d-1}$的线性组合。这会使$\Q[u]$的维数至多$d<4$，与含四个线性无关元素矛盾。因此$f$是次数最小的首一关系，即极小多项式。

\end{proof}

\paragraph{复习提示}

找到一个四次关系只证明极小次数不超过四，还必须排除低次关系。



\subsection{习题36：三次代数关系下的乘法和求逆}

\label{ex:3:36}

\paragraph{题目}

在复数域内设$u$为$f(x)=x^3-x+1$的一根，试将元素
\[
(5u^2+3u-1)(2u^2-2u+6)\quad\text{和}\quad(3u^2-u+2)^{-1}
\]
表成次数不超过$2$的$u$的多项式。

\paragraph{考点与思路}

将$u^3,u^4$化为低次幂，并逐项展开验证逆元。

\paragraph{所用结论与原文依据}

\noteref{c3poly}；\noteref{c3minimal}。原题：《代数学引论》第三版，书页134，扫描PDF第144页。

% 原题逐页视觉核对；扫描PDF第144页；本题证明为本册解答。

\begin{proof}

原题在复数域内计算，所以$7$非零且可逆。由$u^3-u+1=0$，得$u^3=u-1$、$u^4=u^2-u$。

首先展开乘积：
\begin{align*}
(5u^2+3u-1)(2u^2-2u+6)
&=10u^4-4u^3+22u^2+20u-6\\
&=10(u^2-u)-4(u-1)+22u^2+20u-6\\
&=32u^2+6u-2.
\end{align*}
对求逆，取候选$(-2u^2-u+3)/7$。直接展开，
\begin{align*}
(3u^2-u+2)(-2u^2-u+3)
&=-6u^4-u^3+6u^2-5u+6\\
&=-6(u^2-u)-(u-1)+6u^2-5u+6\\
&=7.
\end{align*}
这些因子都是$u$的有理系数多项式，彼此交换，所以反向乘积也为$7$。除以$7$得到
\[
(3u^2-u+2)^{-1}=\frac{-2u^2-u+3}{7}.
\]

\end{proof}

\paragraph{复习提示}

将高次幂逐项约去；求逆的结果须与原元素相乘，核对乘积等于1。



\subsection{习题37：八元域的商环构造}

\label{ex:3:37}

\paragraph{题目}

设$F_2=\Z/(2)$。证明：$f(x)=x^3+x^2+1$在$F_2[x]$内不可约，并且证明商环$R=F_2[x]/(x^3+x^2+1)$是$8$个元素的域。

\paragraph{考点与思路}

先在两个元素上检验三次式无根，再以带余除法数出商环的全部代表元。

\paragraph{所用结论与原文依据}

\noteref{c3irred}；\noteref{c3division}。原题：《代数学引论》第三版，书页134，扫描PDF第144页。

% 原题逐页视觉核对；扫描PDF第144页；本题证明为本册解答。

\begin{proof}

令$h=x^3+x^2+1\in\F_2[x]$。有$h(0)=1$、$h(1)=1+1+1=1$，所以$h$在$\F_2$中无根。
若域上的三次多项式可约，则可分成两个正次数因子，次数和为$3$，必有一个因子的次数为$1$；一次因子的根也是原多项式的根。因此三次多项式无根即不可约，$h$不可约。

由不可约多项式的商环判据，$\F_2[x]/(h)$为域。用$h$作带余除法，每个剩余类有代表元$a+bx+cx^2$，$a,b,c\in\F_2$。若两个这样的代表元同余，其差是次数小于$3$的$h$的倍式，只能为零，所以表示唯一。三系数独立各取两值，因此恰有$2^3=8$个元素。

\end{proof}

\paragraph{复习提示}

无根等价于不可约仅适用于二次或三次式，不能推广到任意次数。



\subsection{习题38：五元域上的全部首一不可约二次式}

\label{ex:3:38}

\paragraph{题目}

设$F_5=\Z/(5)$，试定出$F_5[x]$中全部首项系数为$1$的$2$次不可约多项式。

\paragraph{考点与思路}

保留判别式路线，以表格逐一列出$b$与两种非平方判别式对应的$c$。

\paragraph{所用结论与原文依据}

\noteref{c3irred}；\noteref{c3roots}。原题：《代数学引论》第三版，书页134，扫描PDF第144页。

% 原题逐页视觉核对；扫描PDF第144页；本题证明为本册解答。

\begin{proof}

考虑$f=x^2+bx+c\in\F_5[x]$。域上的二次式可约当且仅当有根，而
\[
4f(x)=(2x+b)^2-(b^2-4c).
\]
因$2$可逆，映射$x\mapsto2x+b$为双射，所以$f$有根当且仅当$\Delta=b^2-4c$是平方（含$0$）。$\F_5$的平方为$0,1,4$，非平方恰为$2,3$。

模$5$下$4=-1$，故$\Delta=b^2+c$，即$c=\Delta-b^2$。计算如下：
\begin{center}\begin{tabular}{c|c|cc}
$b$&$b^2$&$\Delta=2$时的$c$&$\Delta=3$时的$c$\\\hline
0&0&2&3\\1&1&1&2\\2&4&3&4\\3&4&3&4\\4&1&1&2
\end{tabular}\end{center}
因此全部所求多项式为
\begin{align*}
&x^2+2,\quad x^2+3,\\
&x^2+x+1,\quad x^2+x+2,\\
&x^2+2x+3,\quad x^2+2x+4,\\
&x^2+3x+3,\quad x^2+3x+4,\\
&x^2+4x+1,\quad x^2+4x+2.
\end{align*}
表格遍历了全部五个$b$及全部非平方判别式，每对$b,c$唯一决定一个多项式，故无遗漏或重复。

\end{proof}

\paragraph{复习提示}

判别式为零也可约；必须取非平方$2,3$，而不是仅要求判别式非零。



\subsection{习题39：二十五元域的构造}

\label{ex:3:39}

\paragraph{题目}

造一个$25$个元素的域。

\paragraph{考点与思路}

直接使用上一题的不可约二次式，并给出商环中的运算和唯一表示。

\paragraph{所用结论与原文依据}

\noteref{c3irred}；\noteref{c3division}；\exref{3}{38}。原题：《代数学引论》第三版，书页134，扫描PDF第144页。

% 原题逐页视觉核对；扫描PDF第144页；本题证明为本册解答。

\begin{proof}

由\exref{3}{38}，$x^2+2$在$\F_5[x]$中不可约。令
\[
F=\F_5[x]/(x^2+2),\qquad \alpha=x+(x^2+2).
\]
不可约判据保证$F$是域。带余除法说明每个元素唯一写为$a+b\alpha$，其中$a,b\in\F_5$，因此共有$5^2=25$个元素。

其运算由$\alpha^2=-2=3$确定：
\begin{align*}
(a+b\alpha)+(c+d\alpha)&=(a+c)+(b+d)\alpha,\\
(a+b\alpha)(c+d\alpha)&=(ac+3bd)+(ad+bc)\alpha.
\end{align*}
为核对非零元可逆，若$(a,b)\ne(0,0)$，则
\[
(a+b\alpha)(a-b\alpha)=a^2+2b^2\ne0.
\]
若$b\ne0$而右端为零，则$(a/b)^2=-2=3$，与$3$非平方矛盾；若$b=0$，则$a\ne0$也使右端非零。因此逆元为$(a-b\alpha)/(a^2+2b^2)$。

\end{proof}

\paragraph{复习提示}

不能用$\Z/(25)$代替该域；后者有非零零因子$\bar5$。



\subsection{习题40：一个非主理想}

\label{ex:3:40}

\paragraph{题目}

证明$\Z[x]$的理想$(3,x^3+2x^2+2x-1)$不是主理想。

\paragraph{考点与思路}

若该理想由单个多项式生成，该生成元必须整除常数3，因而只能是常数；逐一排除可能的常数。

\paragraph{所用结论与原文依据}

\noteref{c3principal}；\noteref{c3degree}；\noteref{c3polyquot}。原题：《代数学引论》第三版，书页134，扫描PDF第144页。

% 原题逐页视觉核对；扫描PDF第144页；本题证明为本册解答。

\begin{proof}

令$f=x^3+2x^2+2x-1$、$I=(3,f)\subseteq\Z[x]$。先证明$I$是真理想。若$1\in I$，存在$u,v\in\Z[x]$使$1=3u+fv$。逐系数模$3$约化，得到
\[
1=(x^3-x^2-x-1)\bar v\quad\text{于 }\F_3[x].
\]
若$\bar v=0$，右侧为零；若$\bar v\ne0$，右侧次数为$3+\deg\bar v\ge3$。两者都不可能等于$1$，所以$I\ne\Z[x]$。

反设$I=(g)$。由于$3\in I$，有$3=gh$对某个$h\in\Z[x]$成立。两个因子都非零，而$\Z$为整环，故
\[
0=\deg3=\deg g+\deg h.
\]
所以$g,h$均为非零整数常数，且$g\mid3$，即$g=\pm1$或$\pm3$。

若$g=\pm1$，则$(g)=\Z[x]$，与$I$真矛盾。若$g=\pm3$，则$f\in(g)$要求$f$的每个系数都被$3$整除，但$f$的首项系数为$1$，矛盾。因此$I$不能由一个元素生成，不是主理想。

\end{proof}

\paragraph{复习提示}

本题要求非主理想；它是不是极大理想是下一题讨论的另一项性质。



\subsection{习题41：系数约化诱导的商环同构与极大性}

\label{ex:3:41}

\paragraph{题目}

设$F_3=\Z/(3)$，$I=(3,x^3+2x^2+2x-1)$为$\Z[x]$的理想。证明
\[
\Z[x]/I\cong F_3[x]/(x^3-x^2-x-1).
\]
由此证明$I$为$\Z[x]$的一个极大理想。

\paragraph{考点与思路}

先证明复合同态的核恰为题中理想，再独立检验模3三次式不可约。

\paragraph{所用结论与原文依据}

\noteref{c3polyquot}；\noteref{c3irred}；\noteref{c3prime}。原题：《代数学引论》第三版，书页134，扫描PDF第144页。

% 原题逐页视觉核对；扫描PDF第144页；本题证明为本册解答。

\begin{proof}

记$f=x^3+2x^2+2x-1\in\Z[x]$、$h=x^3-x^2-x-1\in F_3[x]$。令$\rho:\Z[x]\to F_3[x]$逐系数模$3$约化，复合取商映射得到满同态
\[
\psi:\Z[x]\longrightarrow F_3[x]/(h),\qquad
g\longmapsto\bar g+(h).
\]
因为$\rho(3)=0$且$\rho(f)=h$，有$I=(3,f)\subseteq\ker\psi$。反过来，若$\psi(g)=0$，则$\bar g=h\bar q$。选取$\bar q$的整数系数提升$q$，$g-fq$的每个系数都被$3$整除，可写$g=fq+3r$，所以$g\in I$。因此$\ker\psi=I$。

由同态基本定理，
\[
\Z[x]/I\cong F_3[x]/(h),\qquad
g+I\longmapsto\bar g+(h).
\]
为证明$I$极大，计算
\[
h(0)=2,\qquad h(1)=1,\qquad h(2)=1\quad\text{于 }F_3.
\]
故$h$没有根。若域上的三次式可约，两个正次数因子中必有一个是一次式，从而会有根；因此$h$不可约。由不可约多项式的商域判据，$F_3[x]/(h)$为域，所以上述同构左侧为域，$I$为极大理想。

\end{proof}

\paragraph{复习提示}

本题的极大性由不可约性推出，不能把上一题的“不是主理想”当成极大性的依据。



\subsection{习题42：商环无非零幂零元与无重因式}

\label{ex:3:42}

\paragraph{题目}

设$F$为一域，$f(x)\in F[x]$，$\deg f(x)>0$。证明商环$F[x]/(f(x))$无幂零元的充要条件是$f(x)$无重因式。

\paragraph{考点与思路}

按不可约因子的指数分析整除性；有重因式时明确构造一个非零幂零剩余类。

\paragraph{所用结论与原文依据}

\noteref{c3division}；\noteref{c3irred}；\noteref{c3degree}；\noteref{c3nil}。原题：《代数学引论》第三版，书页134，扫描PDF第144页。

% 原题逐页视觉核对；扫描PDF第144页；本题证明为本册解答。

\begin{proof}

原题“无幂零元”是指没有非零幂零元，零元本身的幂当然为零。“$f$无重因式”指不存在正次数多项式$p$使$p^2\mid f$。

先说明本题所用的多项式因子性质。$F[x]$中每个正次数多项式都能按次数归纳分解为不可约因子的乘积：可约时分成两个次数更小的正次数因子，继续分解，次数下降保证过程终止。若不可约$p\mid uv$而$p\nmid u$，则$p,u$互素；带余除法的欧几里得算法给出$s,t\in F[x]$使$sp+tu=1$，乘以$v$可得$p\mid v$。故不可约因子有素因子性质，逐次消去还给出分解的唯一性（忽略非零常数和因子顺序）。

设$f$无重因式，写成
\[
f=c\,p_1p_2\cdots p_s,
\]
其中$c\in F^\times$，各$p_i$为互不相同的首一不可约多项式。若$\bar g=g+(f)$在$F[x]/(f)$中幂零，则$g^m\in(f)$对某个$m\ge1$成立，即$f\mid g^m$。每个$p_i\mid g^m$，反复使用素因子性质得$p_i\mid g$。这些互不相同的不可约因子因而都出现在$g$的分解中，所以$p_1\cdots p_s\mid g$，即$f\mid g$。因此$\bar g=0$，商环没有非零幂零元。

反过来，设$f$有重因式。选一个重复的首一不可约因子$p$，写
\[
f=p^e h,\qquad e\ge2,\quad p\nmid h.
\]
令$g=p^{e-1}h=f/p$。它非零，且$\deg g<\deg f$，所以$g\notin(f)$，即$\bar g\ne0$。但
\[
g^2=p^{2e-2}h^2=f\,p^{e-2}h\in(f),
\]
故$\bar g^2=0$。商环出现非零幂零元。两方向合并，得到原题的等价性。

\end{proof}

\paragraph{复习提示}

“无重因式”不同于“不可约”；本题只排除非零幂零元，不要求商环无零因子。



\subsection{习题43：幂零多项式的每个系数均幂零}

\label{ex:3:43}

\paragraph{题目}

设$R$为一个交换环，$f(x)\in R[x]$是一个零幂元。证明$f(x)$的每个系数也都是$R$的零幂元。

\paragraph{考点与思路}

先取最高系数，再模去其生成的幂零理想，以次数归纳处理其余系数。

\paragraph{所用结论与原文依据}

\noteref{c3polyquot}；\noteref{c3nil}；\noteref{c3nilideal}。原题：《代数学引论》第三版，书页134，扫描PDF第144页。

% 原题逐页视觉核对；扫描PDF第144页；本题证明为本册解答。

\begin{proof}

对$f$的次数归纳。若$f=0$，所有系数为零；若$f$为常数，结论就是假设本身。设$f=\sum_{i=0}^m a_ix^i$，$a_m\ne0$、$m\ge1$，且$f^N=0$。$f^N$的$x^{mN}$系数为$a_m^N$，故$a_m^N=0$。

令$I=(a_m)$。因环交换，每个$I$中元素都为$ra_m$，任意$N$个这类元素相乘均为零，所以$I^N=0$。在商环$(R/I)[x]$中，$f$的像$\bar f$仍满足$\bar f^{\,N}=0$，而最高项已经消失。若$\bar f=0$，每个系数都属于$I$；否则其次数至多为$m-1$，由归纳假设，其每个系数均幂零。因此两种情形中，对每个$i<m$都存在$k_i\ge1$使
\[
a_i^{k_i}\in I.
\]
由于$I^N=0$，有$(a_i^{k_i})^N=a_i^{k_iN}=0$，故$a_i$幂零。连同$a_m$的幂零性，归纳完成。

\end{proof}

\paragraph{复习提示}

商环中幂零只给出$a_i^{k_i}\in I$，须利用$I^N=0$才能提升回母环的幂零性。



\subsection{习题44：多项式零因子的常数湮灭元判据}

\label{ex:3:44}

\paragraph{题目}

设$R$为一交换环。证明：若$f(x)\in R[x]$是一个零因子，则存在$R$个元素$a$（$a\ne0$）使得$af(x)=0$；反之显然。

\paragraph{考点与思路}

取最低次的非零湮灭多项式，逐一消去$f$的最高系数，再取一个非零系数作常数湮灭元。

\paragraph{所用结论与原文依据}

\noteref{zerodiv}；\noteref{c3poly}；\noteref{ringlaws}。原题：《代数学引论》第三版，书页134，扫描PDF第144页。

% 原题逐页视觉核对；扫描PDF第144页；本题证明为本册解答。

\begin{proof}

由$f$为零因子，已知$f\ne0$。原题的逆向按非零多项式理解；若把$f=0$也代入，右侧在非零环中成立而左侧不成立，这一边界在此说明，不改写原题。充分性为：若$0\ne a\in R$且$af=0$，将$a$视为非零常数多项式，即得到$f$与非零元素的零乘积，故$f$为零因子。

再证必要性。设$f=\sum_{i=0}^m a_ix^i\ne0$，$a_m\ne0$。因$f$为零因子，存在非零多项式$g$使$fg=0$。在全部这类$g$中选次数$d$最小者，写$g=\sum_{j=0}^d b_jx^j$，$b_d\ne0$。

先由$fg=0$的$x^{m+d}$系数得$a_mb_d=0$。因此$a_mg$若非零，其次数小于$d$；而交换性给出$f(a_mg)=a_m(fg)=0$，与$d$最小矛盾。所以$a_mg=0$，从而$(f-a_mx^m)g=0$。

现作向下归纳。若已知$a_{k+1}g=\cdots=a_mg=0$，则
\[
\left(\sum_{i=0}^k a_ix^i\right)g=0.
\]
比较$x^{k+d}$系数得$a_kb_d=0$，故$a_kg$若非零则次数小于$d$。同时$f(a_kg)=a_k(fg)=0$，仍违背最小性。因此$a_kg=0$。由归纳，对每个$i$都有$a_ig=0$。

比较这些等式的$x^d$系数，得$a_ib_d=0$对全部$i$成立。取非零常数$a=b_d$，便有$af=0$。

\end{proof}

\paragraph{复习提示}

以最低次湮灭多项式作下降论证；逆向仍须按原笔记的非零零因子约定理解。



\subsection{习题45：常数项可逆且其余系数幂零的多项式可逆}

\label{ex:3:45}

\paragraph{题目}

设$R$为一交换环，$f(x)=a_0x^n+a_1x^{n-1}+\cdots+a_n\in R[x]$，若$a_n$为单位而$a_0,a_1,\ldots,a_{n-1}$为零幂元，证明$f(x)$为$R[x]$的单位。

\paragraph{考点与思路}

除去可逆常数项，把剩余多项式证明为幂零，再用有限等比和给出逆。

\paragraph{所用结论与原文依据}

\noteref{unit}；\noteref{c3nil}；\noteref{ringlaws}。原题：《代数学引论》第三版，书页134，扫描PDF第144页。

% 原题逐页视觉核对；扫描PDF第144页；本题证明为本册解答。

\begin{proof}

设$f=a_0x^n+\cdots+a_{n-1}x+a_n$，$a_n$为单位。令
\[
h=a_n^{-1}(f-a_n)=\sum_{i=0}^{n-1}(a_n^{-1}a_i)x^{n-i}.
\]
每个$a_i$幂零，取$m_i\ge1$使$a_i^{m_i}=0$。交换性给出$(a_n^{-1}a_i)^{m_i}=0$。若$n=0$，$f=a_n$已经是单位；以下设$n\ge1$，并令$N=1+\sum_{i=0}^{n-1}(m_i-1)$。

将$h^N$按分配律展开。每项选取$N$个系数因子；若第$i$种系数均只出现至多$m_i-1$次，总数将小于$N$。所以每项至少含某个$(a_n^{-1}a_i)^{m_i}$，从而该项为零。因此$h^N=0$。

令$S=\sum_{j=0}^{N-1}(-h)^j$。有限和相乘消去中间项，得到
\[
(1+h)S=S(1+h)=1-(-h)^N=1.
\]
由于$f=a_n(1+h)$，故
\[
\boxed{f^{-1}=a_n^{-1}\sum_{j=0}^{N-1}(-h)^j.}
\]

\end{proof}

\paragraph{复习提示}

有限多个幂零系数保证存在一个共同的展开次数；不能只说“每个系数幂零，所以和幂零”而不说明理由。



\subsection{习题46：各系数幂零推出多项式幂零}

\label{ex:3:46}

\paragraph{题目}

证明习题43的逆也成立。

\paragraph{考点与思路}

给出统一指数，并明确展开中使用的抽屉计数。

\paragraph{所用结论与原文依据}

\noteref{c3nil}；\exref{3}{43}。原题：《代数学引论》第三版，书页134，扫描PDF第144页。

% 原题逐页视觉核对；扫描PDF第144页；本题证明为本册解答。

\begin{proof}

设$f=\sum_{i=0}^m a_ix^i$，并对每个$i$选$n_i\ge1$使$a_i^{n_i}=0$。令
\[
N=1+\sum_{i=0}^m(n_i-1).
\]
将$f^N$按分配律展开，每项都是$N$个$a_ix^i$的乘积。设第$i$类因子出现$k_i$次，则$\sum_i k_i=N$。若所有$k_i<n_i$，则$\sum_i k_i\le\sum_i(n_i-1)=N-1$，矛盾。因此至少有一个$i$满足$k_i\ge n_i$。

环交换，故该项的系数可重排为$\prod_i a_i^{k_i}$，其中含某个为零的因子$a_i^{n_i}$。所以每项为零，$f^N=0$。连同\exref{3}{43}，得到$f$幂零当且仅当其每个系数幂零。

\end{proof}

\paragraph{复习提示}

指数$N$只需保证充分大，不必是$f$的最小幂零指数。



\subsection{习题47：多元多项式环的幂零根}

\label{ex:3:47}

\paragraph{题目}

设$R$为交换环，$R'=R[x_1,\ldots,x_n]$，用$r(R)$表示$R$的诣零根。证明$R'$的诣零根$r(R')=r(R)[x_1,\ldots,x_n]$。

\paragraph{考点与思路}

逐个添加变量，反复使用一元系数判据。

\paragraph{所用结论与原文依据}

\noteref{c3nil}；\noteref{c3multi}；\exref{3}{43}；\exref{3}{46}。原题：《代数学引论》第三版，书页134，扫描PDF第144页。

% 原题逐页视觉核对；扫描PDF第144页；本题证明为本册解答。

\begin{proof}

依原题，$R$为交换环，$r(R)$表示幂零元全体，也等于全部素理想之交。由\exref{3}{43}与\exref{3}{46}，对$f\in R[x]$，有
\[
f\in r(R[x])\ \Longleftrightarrow\ f\text{的每个系数属于 }r(R).
\]
所以$r(R[x])=r(R)[x]$，右侧指所有系数均属于$r(R)$的多项式集合。

对变量个数$n$归纳。$n=1$已证，$n=0$时也是恒等式。若对$n-1$成立，令$S=R[x_1,\ldots,x_{n-1}]$。它仍是交换幺环，故
\begin{align*}
r(R[x_1,\ldots,x_n])
&=r(S[x_n])\\
&=r(S)[x_n]\\
&=\bigl(r(R)[x_1,\ldots,x_{n-1}]\bigr)[x_n]\\
&=r(R)[x_1,\ldots,x_n].
\end{align*}
最后一步是把关于$x_n$的各个系数再次按前$n-1$个变量展开；最终每个标量系数都在$r(R)$中，反向也成立。

\end{proof}

\paragraph{复习提示}

右侧是系数属于幂零根的多项式理想，不是宣称整个理想有统一幂零指数。



\subsection{习题48：有限环中右逆必为左逆}

\label{ex:3:48}

\paragraph{题目}

在一个有限环$R$中，若一个元素有右逆，则它必有左逆。因而它有逆。当$R$无限时，这个事实对不对？

\paragraph{考点与思路}

保留有限集上乘法映射的证明，并把无限维空间上的移位算子逐个定义。

\paragraph{所用结论与原文依据}

\noteref{unit}；\exref{3}{2}；\exref{3}{3}。原题：《代数学引论》第三版，书页134，扫描PDF第144页。

% 原题逐页视觉核对；扫描PDF第144页；本题证明为本册解答。

\begin{proof}

设幺环$R$有限，且$uv=1$。对每个$y\in R$，有$u(vy)=y$，所以$L_u:x\mapsto ux$满射。有限集合上的满射必单射。又
\[
u(vu)=(uv)u=u=u1,
\]
单射性给出$vu=1$。因此右逆也是左逆。

无限环中不成立。取域$F$上以$e_0,e_1,e_2,\ldots$为基的向量空间$V$（每个向量为有限线性组合），令$R=\operatorname{End}_F(V)$，乘法取映射复合。定义线性映射
\[
v(e_i)=e_{i+1}\ (i\ge0),\qquad u(e_0)=0,\quad u(e_{i+1})=e_i\ (i\ge0).
\]
给定基上的像唯一决定线性映射。对每个$i\ge0$，$(uv)(e_i)=e_i$，所以$uv=\mathrm{id}_V$。但$(vu)(e_0)=0\ne e_0$，故$vu\ne\mathrm{id}_V$。$R$是无限幺环，这给出所需反例。

\end{proof}

\paragraph{复习提示}

无限维向量空间用有限支撑的向量即可定义移位，不需要引入无限和。



\subsection{习题49：多项式环的强拟正则元都是幂零元}

\label{ex:3:49}

\paragraph{题目}

设$R$为一交换环。证明多项式环$R[x]$的强拟正则元都是零幂元。

\paragraph{考点与思路}

取$1-xf$为单位，证明其正次数系数幂零，再用系数判据得到$f$幂零。

\paragraph{所用结论与原文依据}

\noteref{c3nil}；\noteref{c3prime}；\noteref{c3degree}；\exref{3}{22}；\exref{3}{45}；\exref{3}{46}。原题：《代数学引论》第三版，书页134，扫描PDF第144页。

% 原题逐页视觉核对；扫描PDF第144页；本题证明为本册解答。

\begin{proof}

先证所需的一般事实：若$u=\sum a_ix^i$是交换幺环$R$上多项式环的单位，则每个$i>0$的$a_i$幂零。对任一素理想$P$，逐系数取模把$u$送为整环$(R/P)[x]$中的单位。整环上非零乘积次数可加，所以该单位只能是常数，因此$a_i\in P$对所有$i>0$成立。

因$P$任意，这些系数属于全部素理想之交$r(R)$，即均幂零。若$R$为零环，系数全为零，此结论也成立。此外，单位的常数项必为$R$的单位，因为把逆元等式取常数项即可。连同\exref{3}{45}，这也补齐了多项式单位的双向系数判据。

现设$f$在$R[x]$中强拟正则。定义要求对每个$g\in R[x]$，$1-fg$可逆，特别取$g=x$，则$1-xf$可逆。若$f=\sum a_ix^i$，则$1-xf$的$x^{i+1}$系数为$-a_i$，由刚证事实，每个$a_i$都幂零。根据\exref{3}{46}，$f$幂零。

\end{proof}

\paragraph{复习提示}

单位多项式的必要系数条件不能只引用第45题；第45题只证明了充分性，此处已补证必要性。



\subsection{习题50：每个元素满足某个幂等关系时素理想极大}

\label{ex:3:50}

\paragraph{题目}

设$R$为一交换环。若$R$的每个元素$a$都有一个适当方幂$a^n=a$（$n>1$），证明$R$的每个素理想都是极大理想。

\paragraph{考点与思路}

在素理想商整环内使用幂关系，并以无零因子性获得逆元。

\paragraph{所用结论与原文依据}

\noteref{c3prime}；\noteref{field}。原题：《代数学引论》第三版，书页134，扫描PDF第144页。

% 原题逐页视觉核对；扫描PDF第144页；本题证明为本册解答。

\begin{proof}

设$P$为$R$的素理想，$D=R/P$为非零整环。任取$0\ne\bar a\in D$，选代表元$a\in R$。题设给出某个$n>1$使$a^n=a$，故
\[
\bar a^n=\bar a,\qquad \bar a(\bar a^{n-1}-1)=0.
\]
因$\bar a\ne0$且$D$无零因子，得$\bar a^{n-1}=1$。于是$\bar a^{n-2}$是$\bar a$的逆元；当$n=2$时按零次幂为$1$理解。
因此$D$每个非零元均可逆，是域，故$P$为极大理想。指数$n$允许随$a$变化，证明只对当前元素选取其对应的$n$。

\end{proof}

\paragraph{复习提示}

不能在原环中直接消去$a$，因为原环可以有零因子；应先到素理想的商整环。



\subsection{习题51：布尔环的特征、素商与有限生成理想}

\label{ex:3:51}

\paragraph{题目}

如果一个交换环$R$的每个元素$a$满足$a^2=a$，那$R$叫做一个布尔（Boole）环。证明布尔环$R$有如下性质：
\begin{enumerate}[label=(\roman*)]
\item $2a=0$，对所有的$a\in R$；
\item 每个素理想$I$都极大而且$R/I$是特征$2$的素域；
\item $R$的每个有限生成的理想都可由一个元素生成。
\end{enumerate}

\paragraph{考点与思路}

由$(a+a)^2$得到$2a=0$，再展开$(a+b)^2$得到交换性，最后合并两个理想生成元。

\paragraph{所用结论与原文依据}

\noteref{c3prime}；\noteref{c3principal}；\noteref{ringlaws}。原题：《代数学引论》第三版，书页134，扫描PDF第144页。

% 原题逐页视觉核对；扫描PDF第144页；本题证明为本册解答。

\begin{proof}

\begin{enumerate}
\item 对任意$a\in R$，布尔条件给出$a^2=a$及$(a+a)^2=a+a$。分配律展开后有$4a^2=2a$，即$4a=2a$，在加法群中相减得$2a=0$。特别地，$-a=a$。

对$a+b$再用布尔条件，得$a^2+ab+ba+b^2=a+b$，故$ab+ba=0$。由$-ba=ba$，得到$ab=ba$，所以$R$交换。
\item 在本组素理想题目的约定下，$R$有幺元。设$P$为素理想，则$D=R/P$是非零整环，且每个元素$z$仍满足$z^2=z$。由$z(z-1)=0$，每个$z$只能为$0$或$1$。二者因商环非零而不同，所以$D$恰有两个元素，其加法、乘法分别由$1+1=0$和$1\cdot1=1$确定。因此$D\cong\F_2$，且$P$极大。
\item 任取$a,b\in R$，令$c=a+b+ab$。交换性、$a^2=a,b^2=b$及$2ab=0$给出
\[
ac=a^2+ab+a^2b=a,\qquad bc=ab+b^2+ab^2=b.
\]
所以$a,b\in(c)$，得到$(a,b)\subseteq(c)$。另一方面$c\in(a,b)$，故$(c)\subseteq(a,b)$，于是$(a,b)=(c)$。

对生成元个数归纳：一个生成元的理想已是主理想；若$(a_1,\ldots,a_k)=(d)$，则$(a_1,\ldots,a_{k+1})=(d,a_{k+1})$可按上述公式再合并成一个生成元。零理想为$(0)$，所以全部有限生成理想均为主理想。
\end{enumerate}

\end{proof}

\paragraph{复习提示}

先得到$2a=0$再推出交换性，不能在尚未证明交换时把$(a+b)^2$写成$a^2+2ab+b^2$。



\subsection{习题52：唯一素理想、非单位幂零与根商域的等价性}

\label{ex:3:52}

\paragraph{题目}

设$R$为一交换环，$r(R)$为它的诣零根。证明下列三条等价：
\begin{enumerate}[label=(\roman*)]
\item $R$有一个而且只有一个素理想；
\item $R$的每个元素不是单位便是零幂元；
\item $R/r(R)$是一个域。
\end{enumerate}

\paragraph{考点与思路}

按三个条件循环推导，分别使用极大理想存在性、商域判据和诣零根的交表示。

\paragraph{所用结论与原文依据}

\noteref{c3nil}；\noteref{c3prime}；\noteref{c3maxexists}；\exref{3}{20}。原题：《代数学引论》第三版，书页135，扫描PDF第145页。

% 原题逐页视觉核对；扫描PDF第145页；本题证明为本册解答。

\begin{proof}

将三条件记为(1)、(2)、(3)。以下在非零环情形证明；若按一般幺环约定允许零环，则该退化情形中(2)成立、(1)和(3)不成立，是原题须另行排除的边界。题面仍按原书保留。

\begin{enumerate}
\item[(1)$\Rightarrow$(2)] 设唯一素理想为$P$。由幂零根等于全部素理想之交，有$r(R)=P$。$R\ne0$保证至少有一个极大理想，而每个极大理想均为素理想，所以$P$本身是极大理想，且是唯一的极大理想。

若$a\notin P$但不可逆，由\exref{3}{20}，$a$属于某个极大理想，只能是$P$，矛盾。因此$P$外的元素都是单位；$P=r(R)$内的元素都是幂零元。得到(2)。
\item[(2)$\Rightarrow$(3)] 非零环的单位不可能幂零：若$a^m=0$且$a$可逆，乘以$a^{-m}$会得$1=0$。所以$r(R)$不含$1$，为真理想。取$R/r(R)$中任意非零类$a+r(R)$，则$a\notin r(R)$，不是幂零元。由(2)，$a$为单位，故$a^{-1}+r(R)$为该类的逆。于是$R/r(R)$为域。
\item[(3)$\Rightarrow$(1)] 商环为域使$r(R)$为极大理想，特别也是素理想。对任意素理想$P$，由$r(R)$为全部素理想之交，有$r(R)\subseteq P$。因为$r(R)$极大且$P\ne R$，只能$P=r(R)$。因此素理想唯一。
\end{enumerate}

\end{proof}

\paragraph{复习提示}

最后一步需要先由商域得到$r(R)$极大，而不能从“每个素理想包含根”直接推出相等。



\subsection{习题53：只有有限多个理想的整环是域}

\label{ex:3:53}

\paragraph{题目}

证明：如果一个整环只有有限多个理想，则它是一个域。

\paragraph{考点与思路}

保留主理想幂链的路线，先从有限性推出相邻相等，再用整环消去。

\paragraph{所用结论与原文依据}

\noteref{c3principal}；\noteref{zerodiv}；\noteref{field}。原题：《代数学引论》第三版，书页135，扫描PDF第145页。

% 原题逐页视觉核对；扫描PDF第145页；本题证明为本册解答。

\begin{proof}

任取$0\ne a\in R$。考虑下降链
\[
(a)\supseteq(a^2)\supseteq(a^3)\supseteq\cdots.
\]
因$R$只有有限多个理想，这些理想不能全部不同，所以存在$i<j$使$(a^i)=(a^j)$。下降关系给出
\[
(a^i)\supseteq(a^{i+1})\supseteq(a^j)=(a^i),
\]
故$(a^i)=(a^{i+1})$。于是$a^i\in(a^{i+1})$，在交换幺环中可写$a^i=ba^{i+1}$，即
\[
a^i(1-ba)=0.
\]
整环中$a\ne0$使$a^i\ne0$，故$1-ba=0$。交换性给出$ab=ba=1$。因此任一非零元$a$可逆，$R$为域。

\end{proof}

\paragraph{复习提示}

有限多个理想不直接意味着任选某一级就相邻相等，需要利用重复及链的包含关系。



\subsection{习题54：理想集合有极小元的整环是域}

\label{ex:3:54}

\paragraph{题目}

设$R$为一整环。证明：如果$R$的任一由理想组成的非空集合按包含关系都有一个极小元，则$R$为一域。

\paragraph{考点与思路}

在指定的幂主理想集合中取极小元，再展开与上一题相同的消去。

\paragraph{所用结论与原文依据}

\noteref{c3principal}；\exref{3}{53}。原题：《代数学引论》第三版，书页135，扫描PDF第145页。

% 原题逐页视觉核对；扫描PDF第145页；本题证明为本册解答。

\begin{proof}

固定$0\ne a\in R$。集合$\mathcal A=\{(a^m):m\ge1\}$非空，由题设存在其中一个极小元$(a^m)$。因$(a^{m+1})\in\mathcal A$且$(a^{m+1})\subseteq(a^m)$，极小性给出相等。

因此有$b\in R$使$a^m=ba^{m+1}$，即$a^m(1-ba)=0$。$R$为整环且$a\ne0$，所以$a^m\ne0$，得到$ba=1$；交换性使$ab=1$。故每个非零元均可逆，$R$为域。

\end{proof}

\paragraph{复习提示}

“极小元”只须在这一个理想集合内使用，不需要把它当成全部非零理想中的最小理想。



\subsection{习题55：华罗庚恒等式中的可逆性与因子次序}

\label{ex:3:55}

\paragraph{题目}

（华罗庚恒等式）设$a,b$为一环$R$的两个单位，而且$ab-1$也是单位。求证$a-b^{-1}$和$(a-b^{-1})^{-1}-a^{-1}$也是单位，而且满足恒等式
\[
\bigl((a-b^{-1})^{-1}-a^{-1}\bigr)^{-1}=aba-a.
\]
（此题如在第一章已经做过，则可直接应用于下题。）

\paragraph{考点与思路}

先分解$c=(ab-1)b^{-1}$，再将$c^{-1}-a^{-1}$分为三个单位的乘积，反序取逆。

\paragraph{所用结论与原文依据}

\noteref{unit}；\noteref{ringlaws}。原题：《代数学引论》第三版，书页135，扫描PDF第145页。

% 原题逐页视觉核对；扫描PDF第145页；本题证明为本册解答。

\begin{proof}

令$c=a-b^{-1}$。因$a,b,ab-1$均可逆，
\[
c=(ab-1)b^{-1}
\]
是两个单位之积，所以可逆；具体逆元为$bc_1$，其中$c_1=(ab-1)^{-1}$。

接着保持顺序计算
\begin{align*}
c^{-1}(a-c)a^{-1}
&=c^{-1}aa^{-1}-c^{-1}ca^{-1}\\
&=c^{-1}-a^{-1}.
\end{align*}
由于$a-c=b^{-1}$，得到
\[
c^{-1}-a^{-1}=c^{-1}b^{-1}a^{-1}.
\]
右侧为三个单位的乘积，故也是单位。按反序取逆，
\begin{align*}
\bigl((a-b^{-1})^{-1}-a^{-1}\bigr)^{-1}
&=(c^{-1}b^{-1}a^{-1})^{-1}\\
&=abc\\
&=ab(a-b^{-1})=aba-a.
\end{align*}
这就是所要求的华罗庚恒等式，亦可写为$aba=a+((a-b^{-1})^{-1}-a^{-1})^{-1}$。每次取逆的对象都已先证明可逆，计算没有交换因子。

\end{proof}

\paragraph{复习提示}

非交换环中取乘积的逆要颠倒顺序；$abc$不能改成任意其他排列。



\subsection{习题56：共轭不变除子环的中心性}

\label{ex:3:56}

\paragraph{题目}

（嘉当--布饶尔--华（Cartan--Brauer--Hua）定理）设$D$为一体，$C$为它的中心即$C=\{x\in D\mid xg=gx,\text{ 对所有的 }g\in D\}$。又设$R$为$D$的一个子体。证明：如果$dRd^{-1}\subset R$对$D$的所有非零元$d$成立，则$R=D$或者$R\subset C$。（参见华罗庚、万哲先著的《典型群》第一章$\S10$。）

\paragraph{考点与思路}

将第12题中$x,x+1$的比较法用于任意除环，先处理子体外的元素，再处理子体内的元素。

\paragraph{所用结论与原文依据}

\noteref{field}；\noteref{c3anti}；\exref{3}{12}。原题：《代数学引论》第三版，书页135，扫描PDF第145页。

% 原题逐页视觉核对；扫描PDF第145页；本题证明为本册解答。

\begin{proof}

记$C=Z(D)=\{z\in D:zd=dz\text{ 对每个 }d\in D\}$。若$R=D$，结论成立。以下设$R\ne D$，取$x\in D\setminus R$。除子环$R$含$0,1$，所以$x,x+1$都非零。

任取$a\in R$，题设给出$b=xax^{-1}\in R$和$c=(x+1)a(x+1)^{-1}\in R$。有$xa=bx$及$xa+a=cx+c$，故
\[
(b-c)x=c-a.
\]
若$b\ne c$，$R$中存在$(b-c)^{-1}$，从而$x=(b-c)^{-1}(c-a)\in R$，矛盾。故$b=c$，进而$c=a$，于是$xa=ax$。该论证对每个$x\notin R$都成立。

再取$y\in R$，固定一个$x\notin R$。因$x+y\notin R$，它也与$a$交换，所以
\[
ya=(x+y)a-xa=a(x+y)-ax=ay.
\]
因此$a$与$R$内外的每个元素都交换，即$a\in C$。$a$任意，得$R\subseteq C$。

\end{proof}

\paragraph{复习提示}

除子环对非零元素的取逆封闭，是由$b-c\ne0$推出$x\in R$的关键。



\subsection{习题57：实多项式函数环到实数域的同态}

\label{ex:3:57}

\paragraph{题目}

令$R$表示实系数多项式函数全体，按函数的加法和乘法$R$成一环。证明：
\begin{enumerate}[label=(\roman*)]
\item 对每个实数$t$，映射$\eta_t:f(x)\mapsto f(t)$，$f(x)\in R$，是$R$到实数域$\R$的一个同态；
\item $R$到$\R$的任一个同态必是某一个$\eta_t$。
\end{enumerate}

\paragraph{考点与思路}

先验证代入映射，再对一般同态区分零映射与非零映射；非零映射在常数上由保序性确定。

\paragraph{所用结论与原文依据}

\noteref{c3functions1}；\noteref{c3poly}；\noteref{ringiso}。原题：《代数学引论》第三版，书页135，扫描PDF第145页。

% 原题逐页视觉核对；扫描PDF第145页；本题证明为本册解答。

\begin{proof}

将自变量函数记为$s$，$s(c)=c$，把实常数与常函数等同。无限域上的多项式函数定理给出$R=\R[s]\cong\R[x]$。
\begin{enumerate}[label=(\roman*)]
\item 对任意$t\in\R$，代入映射满足
\[
\eta_t(f+g)=f(t)+g(t),\qquad\eta_t(fg)=f(t)g(t).
\]
所以它是环同态；每个$c\in\R$都是常函数$c$的像，故它还满射，并满足$\eta_t(1)=1$。
\item 先交代原题的一处同态约定问题：若采用原笔记“一般环同态不预设保幺”的约定，零映射$R\to\R$也是同态，却不是任何$\eta_t$，因为$\eta_t(1)=1$。所以原题这一项对非零同态成立；全部同态应为零同态以及各$\eta_t$。下面直接从“非零”推出保幺，再完成分类，不在题面中添改条件。

设$\eta$为非零同态。由$\eta(1)^2=\eta(1)$及实数域无零因子，$\eta(1)=0$或$1$。若为$0$，任意$f$都有$\eta(f)=\eta(1f)=\eta(1)\eta(f)=0$，矛盾。因此$\eta(1)=1$。

限制在常数域上，若$c\ne0$，由$\eta(c)\eta(c^{-1})=1$知$\eta(c)\ne0$。若$c>0$，写$c=d^2$且$d\ne0$，便有$\eta(c)=\eta(d)^2>0$。于是$c_1<c_2$时，$\eta(c_2)-\eta(c_1)=\eta(c_2-c_1)>0$，故该限制保序。

加法、乘法及保幺性使$\eta$固定每个有理数。若某个$c\in\R$满足$\eta(c)>c$，选有理数$r$使$c<r<\eta(c)$，保序性却给出$\eta(c)<\eta(r)=r$，矛盾。若$\eta(c)<c$，选$\eta(c)<r<c$，则$\eta(r)<\eta(c)$与$\eta(r)=r$矛盾。因此每个实常数都被固定。

令$t=\eta(s)$。对$f=\sum_i a_is^i$，得
\[
\eta(f)=\sum_i\eta(a_i)\eta(s)^i=\sum_i a_it^i=f(t).
\]
所以$\eta=\eta_t$，而$t=\eta(s)$保证唯一性。若把原题“同态”理解为保幺同态或满同态，这就直接证明原题第二项；一般同态的零映射例外已在上面列出。
\end{enumerate}

\end{proof}

\paragraph{复习提示}

保留原题陈述；一般环同态下的零映射例外必须在解答中说明，不能默默删去。



\subsection{习题58：无限域上非零多项式的非零点集}

\label{ex:3:58}

\paragraph{题目}

设$F$为一个无限域，在多项式环$F[x_1,\ldots,x_n]$中任给两个多项式
\[
f(x_1,\ldots,x_n)\quad\text{和}\quad g(x_1,\ldots,x_n),
\]
而且$g(x_1,\ldots,x_n)\ne0$。证明：如果使$g(c_1,\ldots,c_n)\ne0$的所有点$(c_1,\ldots,c_n)$都使得$f(c_1,\ldots,c_n)=0$，则$f(x_1,\ldots,x_n)=0$。

\paragraph{考点与思路}

先看乘积$fg$的路线，并展开无限域上多元零函数只能来自零多项式的归纳。

\paragraph{所用结论与原文依据}

\noteref{c3functionsN}；\noteref{c3degree}；\noteref{c3roots}。原题：《代数学引论》第三版，书页135，扫描PDF第145页。

% 原题逐页视觉核对；扫描PDF第145页；本题证明为本册解答。

\begin{proof}

设$f,g\in F[x_1,\ldots,x_n]$，$g\ne0$。任取$a\in F^n$。若$g(a)\ne0$，题设给出$f(a)=0$；若$g(a)=0$，乘积值也为零。因此
\[
(fg)(a)=0\qquad(a\in F^n).
\]
下面沿原笔记按变量个数归纳，说明无限域上的多项式若处处取零值，则为零多项式。$n=1$时，非零多项式的根数至多为次数，不可能有无限多个根。

设结论对$n-1$成立，写$h=\sum_{j=0}^m h_j(x_1,\ldots,x_{n-1})x_n^j$，并假设$h$处处为零。固定前$n-1$个坐标，得到关于$x_n$的一元多项式，在无限个点上为零，故各系数值$h_j(a_1,\ldots,a_{n-1})$全为零。因前$n-1$个坐标任意，由归纳假设各$h_j=0$，所以$h=0$。

将该结论用于$h=fg$，得$fg=0$作为多项式恒等式。域上的多元多项式环是整环（逐次将一个变量加到整环上，非零首项系数的乘积保持非零）。因$g\ne0$，得到$f=0$。

\end{proof}

\paragraph{复习提示}

“处处取零值”和“零多项式”的等价性依赖无限域，须先证明再用于乘积。



\subsection{习题59：有限域上原点的指示多项式}

\label{ex:3:59}

\paragraph{题目}

设$F$为$q$个元素的有限域（注意$F$的每个非零元素$a$有$a^{q-1}=1$）。在多项式环$R=F[x_1,\ldots,x_n]$内，多项式$g_0(x_1,\ldots,x_n)=(1-x_1^{q-1})(1-x_2^{q-1})\cdots(1-x_n^{q-1})$具有性质
\[
g_0(a_1,\ldots,a_n)=\begin{cases}
1,&(a_1,\ldots,a_n)=(0,\ldots,0),\\
0,&(a_1,\ldots,a_n)\ne(0,\ldots,0).
\end{cases}\tag{1}
\]

\paragraph{考点与思路}

逐个因子判断，并交代$q-1$次幂公式来自有限乘法群。

\paragraph{所用结论与原文依据}

\noteref{lagrange}；\noteref{c3finite}。原题：《代数学引论》第三版，书页135，扫描PDF第145页。

% 原题逐页视觉核对；扫描PDF第145页；本题证明为本册解答。

\begin{proof}

有限域$F$的非零元素构成$q-1$阶乘法群，所以由拉格朗日定理，对每个$a\ne0$都有$a^{q-1}=1$。又$q\ge2$，故$q-1>0$，$0^{q-1}=0$。

在原点，$g_0$的每个因子$1-x_i^{q-1}$均取$1$，所以
\[
g_0(0,\ldots,0)=1.
\]
在任意非零点$(a_1,\ldots,a_n)$，至少有一个$a_i\ne0$，则其对应因子$1-a_i^{q-1}=0$，整个乘积为零。

因此$g_0$恰在原点取$1$，在$F^n\setminus\{0\}$上取$0$。

\end{proof}

\paragraph{复习提示}

这里$n\ge1$，且有限域至少两个元素，因而所有出现的$q-1$次幂均为正次幂。



\subsection{习题60：原点指示函数代表多项式的次数下界}

\label{ex:3:60}

\paragraph{题目}

$F,R$设如习题59。若多项式$f(x_1,\ldots,x_n)\in R$具有性质
\[
f(a_1,\ldots,a_n)=\begin{cases}
0,&(a_1,\ldots,a_n)=(0,\ldots,0),\\
\text{非零元},&(a_1,\ldots,a_n)\ne(0,\ldots,0),
\end{cases}
\]
则$g(x_1,\ldots,x_n)=1-f(x_1,\ldots,x_n)^{q-1}$具有59题中的$g_0(x_1,\ldots,x_n)$一样具有性质(1)。由此进一步证明$g(x_1,\ldots,x_n)$的次数$\ge n(q-1)$。

\paragraph{考点与思路}

先验证$g$与$g_0$表示同一函数，再用逐变量降次和低次数代表唯一性比较总次数。

\paragraph{所用结论与原文依据}

\noteref{c3functionsN}；\noteref{c3finite}；\exref{3}{59}。原题：《代数学引论》第三版，书页135，扫描PDF第145页。

% 原题逐页视觉核对；扫描PDF第145页；本题证明为本册解答。

\begin{proof}

在原点，$f$取值为零，故$g=1-f^{q-1}$取值为$1$；在其他点，$f$的值非零，故其$q-1$次幂为$1$，$g$取值为零。因此$g$具有\exref{3}{59}中$g_0$的性质(1)。

现将$g$逐变量降次。对任何$e\ge q$，
\[
x_i^e-x_i^{e-q+1}=x_i^{e-q}(x_i^q-x_i).
\]
右侧在$F^n$上处处为零。因此在保持函数值的前提下，每个高于或等于$q$的变量指数都可降低$q-1$，并不增加总次数。逐一处理有限多个单项式和变量，得到$r$，满足每个变量的次数小于$q$，且$\deg r\le\deg g$。

为说明低次数代表唯一，设$h$对每个变量的次数均小于$q$且在$F^n$上恒为零。对$n$归纳。$n=1$时，非零多项式次数小于$q$却有$q$个根，与根数定理矛盾。若$n>1$，写
\[
h=\sum_{j=0}^{q-1}h_j(x_1,\ldots,x_{n-1})x_n^j.
\]
固定前$n-1$个坐标，用一元情形得各$h_j$的值都为零；再对$h_j$使用归纳假设，得所有$h_j=0$，故$h=0$。所以同一函数的两个低次数代表必相等。

$r$与$g_0$表示同一函数，而$g_0$对每个变量的次数也小于$q$，因此$r=g_0$。在$g_0$中，单项式$x_1^{q-1}\cdots x_n^{q-1}$的系数为$(-1)^n\ne0$，所以
\[
\boxed{\deg g\ge\deg r=\deg g_0=n(q-1).}
\]
这就是原题要求的次数下界，所有次数均为总次数。

\end{proof}

\paragraph{复习提示}

原题最后要求的是$g$的次数下界；不能把它替换成$f$的次数下界作为题面。



\subsection{习题61：低于变量数次数的多项式还有非零零点}

\label{ex:3:61}

\paragraph{题目}

（阿廷--谢瓦莱（Artin--Chevalley）定理）设$F,R$设如习题59。设$R$中多项式$f(x_1,\ldots,x_n)$的次数$r<n$，$n$为未定元$x_i$的个数，而且假定$f(0,\ldots,0)=0$。证明：$f(x_1,\ldots,x_n)$除$(0,\ldots,0)$外有另外的零点，即存在$(a_1,\ldots,a_n)\ne(0,\ldots,0)$使得$f(a_1,\ldots,a_n)=0$。

\paragraph{考点与思路}

对有限域全体点求和计数，证明幂和消失，并将域中的零等式转化为整数零点数被特征整除。

\paragraph{所用结论与原文依据}

\noteref{c3finite}；\noteref{c3functionsN}；\exref{3}{59}。原题：《代数学引论》第三版，书页135，扫描PDF第145页。

% 原题逐页视觉核对；扫描PDF第145页；本题证明为本册解答。

\begin{proof}

若$f=0$，任取非零向量即为所求零点。以下设$f\ne0$，记$N$为$f$在$F^n$中的零点个数，$p=\operatorname{char}F$。在域$F$中求和，有
\[
N\cdot1_F=\sum_{a\in F^n}\bigl(1-f(a)^{q-1}\bigr),
\]
因为一个求和项在零点取$1$，非零点取$0$。

先计算$S_k=\sum_{t\in F}t^k$，其中$k=0$时求和项按常数函数$1$理解，故$S_0=q\cdot1_F=0$。对$k>0$，取有限乘法群$F^\times$的生成元$\zeta$。若$q-1\nmid k$，则$\zeta^k\ne1$；乘以$\zeta$置换非零元素，给出$S_k=\zeta^k S_k$，故$S_k=0$。若$q-1\mid k$，则每个非零$t$都满足$t^k=1$，所以$S_k=(q-1)\cdot1_F=-1_F$。

对单项式$x_1^{k_1}\cdots x_n^{k_n}$，逐坐标求和的分配律给出
\[
\sum_{a\in F^n}a_1^{k_1}\cdots a_n^{k_n}=\prod_{i=1}^n S_{k_i}.
\]
此乘积可能非零，必须每个$k_i$都是$q-1$的正倍数，特别是总次数至少为$n(q-1)$。但$f^{q-1}$的总次数为$(q-1)\deg f<n(q-1)$，所以其每个单项式在全体点上的和均为零。常数$1$的总和也是$q^n\cdot1_F=0$。故$N\cdot1_F=0$，即整数$N$被$p$整除。

原点已经是零点，所以$N\ge1$。正整数$N$被$p$整除，故$N\ge p\ge2$，必有另一个不同于原点的零点。

\end{proof}

\paragraph{复习提示}

域中得到的等式是$N\cdot1_F=0$，不是整数$N=0$；它说明零点个数被特征整除。



\subsection{习题62：形式幂级数的逆、整环性与商域表示}

\label{ex:3:62}

\paragraph{题目}

设$R$为交换环，$(a_0,a_1,a_2,\ldots)\in R[[x]]$，$a_i\in R$。令$x=(0,1,0,0,\ldots)$。$(a_0,a_1,a_2,\ldots)$可写成$f=\sum\limits_{i=0}^{\infty}a_ix^i$。证明：
\begin{enumerate}[label=(\roman*)]
\item $f$是$R[[x]]$的单位当且仅当$a_0$是$R$的单位；
\item $R[[x]]$为整环当且仅当$R$为整环；
\item 设$R$为整环，$R[[x]]$的商域记成$F\{x\}$，其中$F$为$R$的商域。则$F\{x\}$的元素$g$可表成
\[
g=x^\nu(a_0+a_1x+a_2x^2+\cdots),\qquad a_i\in F,\ a_0\ne0,\ \nu\in\Z.
\]
\end{enumerate}

\paragraph{考点与思路}

递推逆元系数，再以最低非零项检验整环性，最后直接展开任意两个幂级数之商。

\paragraph{所用结论与原文依据}

\noteref{c3fraction}；\noteref{c3poly}；\noteref{unit}。原题：《代数学引论》第三版，书页136，扫描PDF第146页。

% 原题逐页视觉核对；扫描PDF第146页；本题证明为本册解答。

\begin{proof}

先说明形式幂级数的运算。$R[[x]]$由全部序列$(a_0,a_1,\ldots)$组成，写成$\sum_{i\ge0}a_ix^i$；相等是逐系数相等，加法逐系数进行，乘法的$x^m$系数定义为$\sum_{i=0}^m a_ib_{m-i}$。每个系数仅用有限次运算，不涉及解析收敛。交换幺环$R$的运算律逐系数保证$R[[x]]$仍为交换幺环。

\begin{enumerate}
\item 若$fg=1$，比较常数项得$a_0b_0=1$，所以$a_0$可逆。反过来，若$a_0$为单位，递推定义
\[
b_0=a_0^{-1},\qquad
b_m=-a_0^{-1}\sum_{i=1}^m a_i b_{m-i}\quad(m\ge1).
\]
当定义$b_m$时，右侧只涉及已定义的$b_0,\ldots,b_{m-1}$，故此序列存在且唯一。令$g=\sum b_mx^m$，则$fg$的常数项为$1$，其余第$m$项系数为
\[
a_0b_m+\sum_{i=1}^m a_i b_{m-i}=0.
\]
因此$fg=1$；交换性使$gf=1$，故$f$可逆。
\item 若$R$为整环，设非零$f,g$的最低非零项分别为$a_rx^r,b_sx^s$。$fg$中低于$r+s$次的系数全零，第$r+s$次系数恰为$a_rb_s\ne0$，因此$fg\ne0$。常数$1$非零，所以$R[[x]]$为整环。

反过来，常数嵌入$R\hookrightarrow R[[x]]$为单射，保加法、乘法和单位元。若$R[[x]]$是整环，则其常数子环非零且无零因子，故$R$为整环。
\item 设$F$为$R$的商域，逐系数嵌入$R[[x]]\hookrightarrow F[[x]]$。任取非零分式$f/g$，其中$f,g\in R[[x]]\setminus\{0\}$。分离最低非零项，写成$f=x^r u$、$g=x^s v$，其中$u,v\in R[[x]]$的常数项非零。在$F[[x]]$中，这两个常数项都是单位，故由第一项，$u,v$均可逆。于是
\[
\frac fg=x^{r-s}uv^{-1}=x^\nu(a_0+a_1x+a_2x^2+\cdots),
\]
其中$\nu=r-s\in\Z$，$a_i\in F$，且$a_0=u(0)/v(0)\ne0$。

这可在形式洛朗级数域$F((x))$中理解；它由仅有有限个负次幂项的级数组成。非零洛朗级数分离最低次幂后，其余因子常数项非零，可按第一项求逆，所以确实构成域。上述分式嵌入良定义且单射：$f/g=f'/g'$等价于$fg'=f'g$，而逐系数嵌入保此等式；其像为零只可能$f=0$。

因此证明的是$R[[x]]$商域的每个非零元都有上述表示，并非任意$F$系数洛朗级数都必来自该商域。当$R=F$时，任一洛朗级数均可写为$x^{-m}h$，$h\in F[[x]]$，故$\operatorname{Frac}(F[[x]])=F((x))$。
\end{enumerate}

\end{proof}

\paragraph{复习提示}

形式级数不需要收敛；一般$R$时只得到其商域嵌入$F((x))$，不能无条件宣称二者相等。



\subsection{习题63：非单位构成唯一极大理想的条件}

\label{ex:3:63}

\paragraph{题目}

证明：如果在交换环$R$中每个非单位$a$与$1$的和$1+a$为单位，则$R$的全部非单位构成$R$的一个极大理想，而且是$R$的唯一极大理想。这种环存在否？

\paragraph{考点与思路}

把假设的单位和归一化为$1$，用$1+a$的条件排除两个非单位之和为单位。

\paragraph{所用结论与原文依据}

\noteref{c3prime}；\exref{3}{30}。原题：《代数学引论》第三版，书页136，扫描PDF第146页。

% 原题逐页视觉核对；扫描PDF第146页；本题证明为本册解答。

\begin{proof}

以下在$R\ne0$的情形证明；若允许零环，非单位集为空，不构成极大理想，故该退化情形需另行排除。令$M=R\setminus R^\times$，则$0\in M$而$1\notin M$。若$a\in M$、$r\in R$且$ra$是单位，设$(ra)c=1$，交换性给出$a(rc)=1$，使$a$为单位，矛盾。因此$ra\in M$。若$-a$可逆，则$a=(-1)(-a)$也可逆，故$-a\in M$。

任取$a,b\in M$。若$a+b=u$为单位，则$au^{-1},bu^{-1}$均不是单位；否则再乘单位$u$会使$a$或$b$成为单位。于是$-au^{-1}$为非单位，题设给出
\[
1-au^{-1}=bu^{-1}\in R^\times,
\]
矛盾。因此$a+b\in M$。结合上述性质，$M$是理想，且因不含$1$而为真理想。

任一真理想都不能含单位，所以均包含于$M$。若$M\subsetneq J\lhd R$，则$J$含一个不属于$M$的单位，故$J=R$。所以$M$极大；任一极大理想既含于$M$，就只能等于$M$，故唯一。

这种环存在。例如$R=\Z/(p^n)$，$n\ge1$：非单位恰为$p$的倍数，若$p\mid a$，则$p\nmid1+a$，所以$1+\bar a$为单位。唯一极大理想是$(p)/(p^n)$。

\end{proof}

\paragraph{复习提示}

归一化时仍须确认$au^{-1}$非单位；零环需要排除，否则非单位集为空而非理想。



\subsection{习题64：单位加零构成子域及其外元素超越性}

\label{ex:3:64}

\paragraph{题目}

证明：如果在整环$R$中每个非单位$a$（$0$除外）与$1$的和$1+a$仍为非单位，
\begin{enumerate}[label=(\roman*)]
\item $R$的全部单位添加$0$构成$R$的一个子域，记作$F$；
\item $R$中$F$以外的元素都是$F$上的超越元。
\end{enumerate}

\paragraph{考点与思路}

用两个单位之和归一化制造矛盾，再把代数关系化为常数项非零的关系。

\paragraph{所用结论与原文依据}

\noteref{field}；\noteref{c3algebraic}；\noteref{unit}。原题：《代数学引论》第三版，书页136，扫描PDF第146页。

% 原题逐页视觉核对；扫描PDF第146页；本题证明为本册解答。

\begin{proof}

\begin{enumerate}[label=(\roman*)]\item 令$F=R^\times\cup\{0\}$。整环要求$1\ne0$，故$F$含$0,1$。单位的积及逆仍为单位，单位的负元也是单位，所以$F$对乘法、取负和非零元取逆均封闭。

只需检查加法。若一个加数为零，无需再证。设$u,v$均为单位，且$u+v\ne0$。反设$u+v$不是单位，令$a=-(u+v)u^{-1}$。乘以单位保持非单位性，且整环中乘单位保持非零，故$a$是非零非单位。题设要求$1+a$也不是单位，但
\[
1+a=1-(u+v)u^{-1}=-vu^{-1}
\]
是单位，矛盾。因此$u+v$不是零就是单位，$F$对加法封闭，是$R$的子域。

\item 任取$z\in R\setminus F$，则$z\ne0$且不是单位。若$z$在$F$上代数，存在非零多项式$\sum_{i=0}^m c_iX^i$使$\sum c_iz^i=0$。取最小的$c_k\ne0$，因$R$为整环且$z\ne0$，消去公共因子$z^k$，得到常数项非零的关系。重新编号可写为
\[
c_0+c_1z+\cdots+c_dz^d=0,\qquad c_0\ne0.
\]
此处$d\ge1$，否则非零常数为零。因$c_0$是$F$中的非零元，故在$R$中可逆，且
\[
z\bigl[-c_0^{-1}(c_1+c_2z+\cdots+c_dz^{d-1})\bigr]=1.
\]
交换性使括号内元素为$z$的双侧逆，矛盾。因此$z$在$F$上超越。

\end{enumerate}

\end{proof}

\paragraph{复习提示}

若最初关系常数项为零，要利用$z\ne0$和整环消去，不能直接假定有非零常数项。



\subsection{习题65：原题的常数反例与非常数情形}

\label{ex:3:65}

\paragraph{题目}

证明：若$F$为一无限域，则不存在非零多项式$f(x)\in F[x]$使得$f(x)f(y)=f(x+y)$，其中$x,y$在$F$上代数无关。

\paragraph{考点与思路}

先保留并检验原书的“非零”表述；常数1是反例，随后说明排除常数后为何不可能。

\paragraph{所用结论与原文依据}

\noteref{c3degree}；\noteref{c3algebraic}。原题：《代数学引论》第三版，书页136，扫描PDF第146页。

% 原题逐页视觉核对；扫描PDF第146页；本题证明为本册解答。

\begin{proof}

按所核对扫描本的原文，原题写的是“不存在非零多项式”。这一表述不成立：取
\[
f(x)=1,
\]
则$f$为非零多项式，且无论$x,y$是否在$F$上代数无关，都有$f(x)f(y)=1=f(x+y)$。所以原题不能按字面证明。这里保留原题，不把“非零”直接改写进题面为“非常数”。

若所要证明的是不存在非常数解，则结论成立。设$\deg f=d>0$，首项系数为$c\ne0$。因$x,y$在$F$上代数无关，$f(x)f(y)=f(x+y)$等价于$F[X,Y]$中的形式多项式恒等式。左侧$X^dY^d$的系数为$c^2\ne0$，总次数为$2d$；右侧$f(X+Y)$的所有项总次数至多为$d$，故$X^dY^d$系数为零。这与$2d>d$矛盾。

因此全部可能的解只能是常数。设$f=c\in F$，原等式为$c^2=c$，即$c(c-1)=0$。域无零因子，故$c=0$或$c=1$。全部解恰为$0,1$，其中唯一非零解为$1$。

\end{proof}

\paragraph{复习提示}

原书此处存在常数反例；题面照录，辨析和可成立的结论在解答中明确区分。



\subsection{习题66：全矩阵环的左理想都是主左理想}

\label{ex:3:66}

\paragraph{题目}

设$F$为一域，$M_n(F)$为$F$上全矩阵环。证明：$M_n(F)$的每个左理想都是主理想。

\paragraph{考点与思路}

构造左理想对应的行空间，证明每个行向量都能单独放到任意行，再以一组基组成生成矩阵。

\paragraph{所用结论与原文依据}

\noteref{ideal}；\noteref{subring}。原题：《代数学引论》第三版，书页136，扫描PDF第146页。

% 原题逐页视觉核对；扫描PDF第146页；本题证明为本册解答。

\begin{proof}

左理想$L\subseteq M_n(F)$是加法子群，且对任意$X\in M_n(F),A\in L$有$XA\in L$。设$W$为$L$中所有矩阵的所有行向量组成的集合。下面先证它为$F^n$的子空间，而不预先假定这些行都来自同一个矩阵。

若$w$为$A\in L$的第$i$行，则$E_{ki}A\in L$，其第$k$行为$w$、其他行为零。于是任意两个出现过的行$w,z$都能放到同一个指定行后相加；乘以标量矩阵$cE_n$还能把行变为$cw$。零矩阵给出零行，因此$W$为子空间。

记$L_W$为每行都属于$W$的矩阵集合。按定义$L\subseteq L_W$。反过来，若$B\in L_W$，对其每行$w_k$，选一个$L$中矩阵使其某行等于$w_k$，再左乘矩阵单位把它搬到第$k$行、其余行置零。将这$n$个矩阵相加，得到$B\in L$。故$L=L_W$。

取$W$的一组基$w_1,\ldots,w_r$，$0\le r\le n$，令$A$的前$r$行为这些基向量，其余行为零。因$A\in L_W=L$，有$M_n(F)A\subseteq L$。任给$B\in L$，每行$b_i$可写为$\sum_{j=1}^r c_{ij}w_j$。令$X$的前$r$列为这些$c_{ij}$，其余列为零，则$XA=B$。因此$L\subseteq M_n(F)A$，得到
\[
L=M_n(F)A.
\]
当$W=0$时取$A=0$，同样成立。故每个左理想都是主左理想。

\end{proof}

\paragraph{复习提示}

左乘矩阵形成原矩阵各行的线性组合，因此左理想对应行空间；不能在此换成列空间。



\subsection{习题67：阿贝尔群的自同态环}

\label{ex:3:67}

\paragraph{题目}

设$G$为一个阿贝尔群，$\eta$与$\xi$为$G$的同态。我们定义$\eta$与$\xi$的积$\eta\xi$为
\[
(\eta\xi)(a)=\eta(\xi(a)),\qquad a\in G.
\]
显然$G$的自同态全体$\operatorname{End}(G)$形成一个乘法半群。由于$G$的运算交换，我们还可以定义$\eta$与$\xi$的和$\eta+\xi$为
\[
(\eta+\xi)(a)=\eta(a)+\xi(a),\qquad a\in G.
\]
验证$\eta+\xi$是$G$的一个自同态。显然$\operatorname{End}(G)$构成一个环，有单位元素，叫做交换群$G$的自同态环。

\paragraph{考点与思路}

逐项验证两种运算，特别写出依赖自同态加法性的分配律。

\paragraph{所用结论与原文依据}

\noteref{ring}；\noteref{hom}；\noteref{aut}。原题：《代数学引论》第三版，书页136，扫描PDF第146页。

% 原题逐页视觉核对；扫描PDF第146页；本题证明为本册解答。

\begin{proof}

将$G$的群运算写为加法。对$\eta,\xi\in\operatorname{End}(G)$定义
\[
(\eta+\xi)(a)=\eta(a)+\xi(a),\qquad
(\eta\xi)(a)=\eta(\xi(a)).
\]
若$a,b\in G$，因$G$交换，
\begin{align*}
(\eta+\xi)(a+b)
&=\eta(a)+\eta(b)+\xi(a)+\xi(b)\\
&=(\eta(a)+\xi(a))+(\eta(b)+\xi(b)).
\end{align*}
所以$\eta+\xi$仍为自同态。又$(\eta\xi)(a+b)=\eta(\xi(a)+\xi(b))=\eta\xi(a)+\eta\xi(b)$，故乘法也封闭。

逐点加法由$G$的交换加法继承结合律、交换律；零映射为零元，$(-\eta)(a)=-\eta(a)$为加法逆元，且$-\eta$仍为自同态。因此加法构成阿贝尔群。映射复合本来结合，恒等映射$\mathrm{id}_G$为双侧乘法单位。

对任意$a\in G$，有
\begin{align*}
[\eta(\xi+\zeta)](a)
&=\eta(\xi(a)+\zeta(a))=\eta\xi(a)+\eta\zeta(a),\\
[(\eta+\xi)\zeta](a)
&=\eta(\zeta(a))+\xi(\zeta(a))=\eta\zeta(a)+\xi\zeta(a).
\end{align*}
映射相等可逐点判断，故左右分配律成立，$\operatorname{End}(G)$为幺环。当$G=0$时，它是允许$0=1$的零环。

\end{proof}

\paragraph{复习提示}

阿贝尔条件用于保证逐点和仍为群同态；乘法是复合，不是逐点相乘。



\subsection{习题68：整数群、有限循环群、有理数群与初等p群的自同态环}

\label{ex:3:68}

\paragraph{题目}

试决定下列阿贝尔群的自同态环：
\begin{enumerate}[label=(\roman*)]
\item 整数加法群$\Z^+$；
\item $n$阶循环群$G=\langle\sigma\rangle$；
\item 有理数加法群$\Q^+$；
\item $p^n$阶的初等$p$群，即$n$个$p$阶群的直和，$p$为素数。
\end{enumerate}

\paragraph{考点与思路}

依原题四项分别确定自同态，重点区分有限n阶循环群的参数是模n剩余类。

\paragraph{所用结论与原文依据}

\noteref{cyclic}；\noteref{hom}；\exref{3}{67}。原题：《代数学引论》第三版，书页136，扫描PDF第146页。

% 原题逐页视觉核对；扫描PDF第146页；本题证明为本册解答。

\begin{proof}

各环乘法均按\exref{3}{67}取映射复合，右侧映射先作用。
\begin{enumerate}[label=(\roman*)]
\item 对$\Z^+$，任一自同态$\eta$由$m=\eta(1)\in\Z$唯一决定：加法性给出$\eta(k)=km$对每个整数$k$成立，包括负整数。反过来，$\eta_m(k)=mk$是自同态。映射$m\mapsto\eta_m$为双射，且
\[
\eta_m+\eta_n=\eta_{m+n},\qquad\eta_m\eta_n=\eta_{mn},\qquad\eta_1=\mathrm{id}.
\]
因此$\operatorname{End}(\Z^+)\cong\Z$。
\item 设$G=\langle\sigma\rangle$的阶为$n$。任一自同态的生成元像为$\sigma^m$，$m$由模$n$剩余类决定，且
\[
\eta_m(\sigma^k)=\sigma^{mk}.
\]
这个定义良好：若$k\equiv k'\pmod n$，则$n\mid m(k-k')$，两像相等；它保持群乘法。每个自同态都唯一如此给出，所以$\bar m\mapsto\eta_m$是$\Z/(n)$到$\operatorname{End}(G)$的双射。

原题以乘法写群运算，因而逐点相加的含义为$(\eta_m+\eta_\ell)(g)=\eta_m(g)\eta_\ell(g)$。由生成元的像计算，
\[
\eta_m+\eta_\ell=\eta_{m+\ell},\qquad\eta_m\eta_\ell=\eta_{m\ell}.
\]
恒等映射对应$\bar1$，故$\operatorname{End}(G)\cong\Z/(n)$。当$n=1$时，两边均为零环。
\item 对$\Q^+$，令$r=\eta(1)\in\Q$。任取$m\in\Z,n>0$，有
\[
n\eta(m/n)=\eta(m)=mr,
\]
故$\eta(m/n)=r(m/n)$。因此自同态恰为$\mu_r(q)=rq$，$r\in\Q$。反过来每个这样的映射都保加法，且$\mu_r+\mu_s=\mu_{r+s}$、$\mu_r\mu_s=\mu_{rs}$、$\mu_1=\mathrm{id}$。所以$\operatorname{End}(\Q^+)\cong\Q$。
\item 以加法写$G=\bigoplus\limits_{i=1}^n C_p$，每个分量选生成元$e_i$。每个元素唯一为$\sum_i a_ie_i$，$a_i\in\F_p$。由$pg=0$，定义$\bar m g=mg$不依赖整数代表元，于是$G$为$\F_p$上的$n$维向量空间。

任一群自同态满足$T(mg)=mT(g)$，故自动为$\F_p$线性映射；反向也成立。将$T(e_j)=\sum_i t_{ij}e_i$写为列坐标，得到矩阵$A_T=(t_{ij})$。每个$n\times n$矩阵唯一确定这样的$T$，且
\[
A_{T+S}=A_T+A_S,\qquad A_{TS}=A_TA_S,\qquad A_{\mathrm{id}}=E_n.
\]
所以$\operatorname{End}(G)\cong M_n(\F_p)$。
\end{enumerate}

\end{proof}

\paragraph{复习提示}

第二项是有限n阶循环群，自同态环为模n整数环；不能替换成无限循环群。


